%=======================================================================
%  GK_COMPLETE_CALCULATIONS.tex
%
%  COMPLETE MATHEMATICAL AND NUMERICAL CALCULATION RECORD
%  One-dimensional Guyer-Krumhansl heat conduction treated as an
%  inverse problem for (alpha, tau_q, kappa^2).
%
%  THIS IS NOT A JOURNAL MANUSCRIPT.
%  It is a technical calculation document. Its purpose is that another
%  researcher can reconstruct every equation, transformation, parameter
%  value and numerical criterion without relying on hidden reasoning.
%
%  PHASE 1 deliverable of the calculation-first workflow:
%    PHASE 1 calculations -> PHASE 2 bibliography -> PHASE 3 program
%    -> PHASE 4 user runs program -> PHASE 5 compare -> PHASE 6 paper
%
%  Bibliography: references.bib
%  Compile: pdflatex -> bibtex -> pdflatex -> pdflatex
%=======================================================================
\documentclass[11pt,a4paper]{article}

\usepackage[T1]{fontenc}
\usepackage[utf8]{inputenc}
\usepackage{lmodern}
\usepackage[english]{babel}
\usepackage[margin=2.3cm]{geometry}
\usepackage{amsmath,amssymb,amsthm,bm}
\usepackage{graphicx,booktabs,tabularx,longtable,array}
\usepackage{caption}
\usepackage{float}
\usepackage{xcolor}
\usepackage{hyperref}
\hypersetup{colorlinks=true,linkcolor=blue!45!black,citecolor=blue!45!black,
            urlcolor=blue!45!black}
\usepackage{cleveref}
\usepackage{microtype}
\usepackage{siunitx}
\usepackage{enumitem}

%------------------------- structured environments -------------------------
\theoremstyle{definition}
\newtheorem{definitionx}{Definition}[section]
\newtheorem{assumptionx}{Assumption}[section]
\newtheorem{derivationx}{Derivation}[section]
\newtheorem{calculationx}{Calculation}[section]
\newtheorem{propositionx}{Proposition}[section]

\newenvironment{checkbox}%
  {\par\medskip\noindent\begin{minipage}{\linewidth}%
   \rule{\linewidth}{0.4pt}\par\noindent\textbf{CHECK.}\enspace}%
  {\par\rule{\linewidth}{0.4pt}\end{minipage}\par\medskip}

\newenvironment{resultbox}%
  {\par\medskip\noindent\begin{minipage}{\linewidth}%
   \rule{\linewidth}{1.0pt}\par\noindent\textbf{RESULT.}\enspace}%
  {\par\rule{\linewidth}{1.0pt}\end{minipage}\par\medskip}

\newcommand{\source}[1]{\par\smallskip\noindent\textbf{SOURCE.}\enspace #1\par}

%--------------------------- provenance labels ----------------------------
\newcommand{\LIT}{\textcolor{blue!55!black}{\footnotesize[\textsc{literature value}]}}
\newcommand{\DER}{\textcolor{green!35!black}{\footnotesize[\textsc{derived value}]}}
\newcommand{\NUM}{\textcolor{red!55!black}{\footnotesize[\textsc{numerically computed here}]}}
\newcommand{\PREV}{\textcolor{orange!70!black}{\footnotesize[\textsc{previously verified}]}}

%--------------------------- notation macros ------------------------------
\newcommand{\tq}{\tau_{q}}
\newcommand{\kk}{\kappa^{2}}
\newcommand{\Bnum}{B}
\newcommand{\Lap}[1]{\hat{#1}}

\title{\textbf{Complete Calculation Record}\\[0.4ex]
       \large One-Dimensional Guyer--Krumhansl Heat Conduction\\
       as an Inverse Problem for $(\alpha,\tq,\kk)$}
\author{Technical calculation document --- not a manuscript}
\date{23 August 2026}

\begin{document}
\maketitle

\noindent
\rule{\linewidth}{1pt}

\textbf{Purpose and status of this document.}
This document records every calculation underlying the present study in full,
including intermediate algebra that would normally be compressed. It is written
so that the program of Phase~3 can be implemented directly from it, and so that
the program output of Phase~4 can be compared against it line by line in
Phase~5. It contains no introduction, no discussion, no conclusions and no
novelty claims.

\textbf{Provenance labelling.} Every quantitative statement carries one of four
labels, and these are never mixed:
\begin{itemize}[nosep]
  \item \LIT\ a value taken from a published source, cited at the point of use;
  \item \DER\ a value obtained by algebra from literature values or from
        equations derived in this document;
  \item \NUM\ a value computed numerically in this work, with the script named;
  \item \PREV\ a value established in earlier verified work and not recomputed
        in this pass; explicitly flagged as such.
\end{itemize}

\textbf{Reproducibility.} Every numerical result quoted below names the script
that produced it. Scripts live in \texttt{validation/calc/} and
\texttt{repro/}. The audit records are
\texttt{checks\_symbolic.json} (30 symbolic and dimensional checks) and
\texttt{checks\_numerical.json} (19 numerical checks), both in
\texttt{validation/calc/}.

\noindent
\rule{\linewidth}{1pt}

\tableofcontents
\clearpage

%=======================================================================
\section{Problem definition}
\label{sec:problem}
%=======================================================================

\begin{definitionx}[Physical problem]
\label{def:problem}
A homogeneous, rigid, isotropic slab of thickness $L$ is initially in thermal
equilibrium with its surroundings. At $t=0$ the face $x=0$ receives a heat
pulse of finite duration $t_{p}$. The opposite face $x=L$ is thermally
insulated. The measured quantity is the temperature history of the rear face,
$T(L,t)$. This is the laser flash configuration.
\end{definitionx}

\begin{definitionx}[Domain and fields]
\label{def:domain}
\begin{align}
  \text{spatial domain:}&\quad x \in [0,L], \qquad L>0 ,
  \label{eq:domain}\\
  \text{time domain:}&\quad t \geq 0 ,
  \label{eq:time}\\
  \text{temperature field:}&\quad T(x,t) \quad [\si{\kelvin}] ,
  \label{eq:Tfield}\\
  \text{heat flux field:}&\quad q(x,t) \quad
     [\si{\watt\per\metre\squared}] .
  \label{eq:qfield}
\end{align}
The problem is strictly one-dimensional: no quantity depends on $y$ or $z$, and
$q$ denotes the $x$-component of the flux vector.
\end{definitionx}

\begin{definitionx}[Material parameters]
\label{def:params}
The medium is characterised by four constants, all positive:
\begin{align}
  \rho c &\quad \text{volumetric heat capacity}
          \quad [\si{\joule\per\metre\cubed\per\kelvin}] ,
  \label{eq:rhoc}\\
  \lambda &\quad \text{thermal conductivity}
          \quad [\si{\watt\per\metre\per\kelvin}] ,
  \label{eq:lambda}\\
  \tq &\quad \text{relaxation time of the heat flux}
          \quad [\si{\second}] ,
  \label{eq:tauq}\\
  \kk &\quad \text{Guyer--Krumhansl nonlocal coefficient}
          \quad [\si{\metre\squared}] .
  \label{eq:kappa2}
\end{align}
The thermal diffusivity is the derived combination
$\alpha=\lambda/(\rho c)$, in \si{\metre\squared\per\second}.
\end{definitionx}

\begin{assumptionx}[Constitutive assumptions]
\label{ass:const}
\begin{enumerate}[nosep,label=(A\arabic*)]
  \item The solid is rigid: no mechanical deformation, no coupling of the
        temperature field to a displacement field.
  \item $\rho$, $c$, $\lambda$, $\tq$ and $\kk$ are constants, independent of
        temperature and position.
  \item There is no internal heat generation.
  \item The heat flux obeys the Guyer--Krumhansl constitutive law
        \eqref{eq:gk-dim} below, which is linear in $q$, $\partial_{t}q$,
        $\partial_{xx}q$ and $\partial_{x}T$.
  \item All coefficients are non-negative, as required for thermodynamic
        admissibility.
\end{enumerate}
\end{assumptionx}

\begin{definitionx}[Inverse problem]
\label{def:inverse}
The forward problem maps $(\alpha,\tq,\kk)$ to the rear-face history
$T(L,t)$. The inverse problem is the recovery of the parameter vector
\begin{equation}
  \bm{\theta}=(\alpha,\ \tq,\ \kk)
  \label{eq:theta-def}
\end{equation}
from a finite, noisy sample of $T(L,t)$. The central question of the
calculations below is which components or combinations of $\bm{\theta}$ are
determined by that observation.
\end{definitionx}

\subsection{Explicit exclusion list}
\label{sec:exclusions}

The following are \emph{not} part of the model and appear nowhere in this
document or in the program to be built from it. This list is fixed and is not
to be extended later.

\begin{center}
\begin{tabular}{@{}lp{8.6cm}@{}}
\toprule
\textbf{Excluded} & \textbf{Reason} \\
\midrule
Thermoelastic coupling & rigid conductor, Assumption (A1) \\
Piezoelectric coupling & no electric field in the model \\
Two-temperature theory & single temperature field, \eqref{eq:Tfield} \\
Fractional-order derivatives & constitutive law is integer order \\
Memory-dependent derivatives & no memory kernel \\
Rabotnov kernels & no memory kernel \\
Porosity as a mechanism & effective homogeneous medium \\
Semiconductor carrier coupling & no carrier transport equation \\
Klein--Gordon nonlocal terms & not present in \eqref{eq:gk-dim} \\
Nonlocal elasticity (Eringen) & no mechanical field \\
Moore--Gibson--Thompson (MGT) & requires a thermal-displacement variable
  absent from the transfer function \eqref{eq:m2} \\
Multi-dimensional geometry & strictly 1-D, \cref{def:domain} \\
\bottomrule
\end{tabular}
\end{center}

%=======================================================================
\section{Dimensional Guyer--Krumhansl equations}
\label{sec:dimensional}
%=======================================================================

\subsection{Energy balance}

\begin{derivationx}[Balance of internal energy]
\label{der:energy}
For a rigid conductor the specific internal energy is $e=cT$. Conservation of
energy in one dimension, with no source term, reads
$\rho\,\partial_{t}e + \partial_{x}q = 0$. Substituting $e=cT$ with $c$
constant gives
\begin{equation}
  \boxed{\ \rho c\,\frac{\partial T}{\partial t}
        + \frac{\partial q}{\partial x} = 0\ }
  \label{eq:energy-dim}
\end{equation}
\end{derivationx}

\source{Standard balance law for a rigid heat conductor; used in exactly this
form by the validation anchor \cite{Kovacs2018} and by
\cite{Both2016,Van2017}. Not specific to any constitutive assumption.}

\begin{checkbox}
Term-by-term dimensions of \eqref{eq:energy-dim}:
\begin{align*}
  \left[\rho c\,\partial_{t}T\right]
    &= \si{\joule\per\metre\cubed\per\kelvin}\cdot
       \si{\kelvin\per\second}
     = \si{\joule\per\metre\cubed\per\second}
     = \si{\watt\per\metre\cubed} , \\
  \left[\partial_{x}q\right]
    &= \si{\watt\per\metre\squared}\cdot\si{\per\metre}
     = \si{\watt\per\metre\cubed} .
\end{align*}
Both terms are \si{\watt\per\metre\cubed}. Dimensionally consistent.
\end{checkbox}

\subsection{Constitutive law}

\begin{derivationx}[Guyer--Krumhansl constitutive equation]
\label{der:gk}
The Guyer--Krumhansl law extends the Maxwell--Cattaneo--Vernotte relation
$\tq\,\partial_{t}q + q + \lambda\,\partial_{x}T = 0$ by a term proportional to
the second spatial derivative of the flux:
\begin{equation}
  \boxed{\ \tq\,\frac{\partial q}{\partial t}
        + q
        + \lambda\,\frac{\partial T}{\partial x}
        - \kk\,\frac{\partial^{2} q}{\partial x^{2}} = 0\ }
  \label{eq:gk-dim}
\end{equation}
\end{derivationx}

\source{Originally obtained by Guyer and Krumhansl from the linearised phonon
Boltzmann equation for dielectric crystals
\cite{GuyerKrumhansl1966a,GuyerKrumhansl1966b}. An equivalent structure follows
from non-equilibrium thermodynamics with internal variables and entropy-current
multipliers, in which the coefficients are constrained only by the second law
and no carrier mechanism is assumed \cite{Van2001,Nyiri1991}. The flux form
\eqref{eq:gk-dim} is the form used by the validation anchor \cite{Kovacs2018}.}

\begin{checkbox}
Term-by-term dimensions of \eqref{eq:gk-dim}, target
\si{\watt\per\metre\squared}:
\begin{align*}
  \left[\tq\,\partial_{t}q\right]
    &= \si{\second}\cdot\si{\watt\per\metre\squared\per\second}
     = \si{\watt\per\metre\squared} ,\\
  \left[q\right] &= \si{\watt\per\metre\squared} ,\\
  \left[\lambda\,\partial_{x}T\right]
    &= \si{\watt\per\metre\per\kelvin}\cdot\si{\kelvin\per\metre}
     = \si{\watt\per\metre\squared} ,\\
  \left[\kk\,\partial_{xx}q\right]
    &= \si{\metre\squared}\cdot
       \si{\watt\per\metre\squared}\cdot\si{\per\metre\squared}
     = \si{\watt\per\metre\squared} .
\end{align*}
All four terms are \si{\watt\per\metre\squared}. Dimensionally consistent, and
$\kk$ must have dimension \si{\metre\squared} for this to hold; this fixes the
meaning of the symbol.
\end{checkbox}

\subsection{Notation warning: \texorpdfstring{$\kk$}{kappa2} versus a length scale}
\label{sec:notation-warning}

Different sources parameterise the nonlocal term differently, and the
conversions must be stated rather than assumed.

\begin{itemize}[nosep]
  \item This document and the anchor \cite{Kovacs2018} use $\kk$ with dimension
        \si{\metre\squared}, multiplying $\partial_{xx}q$ in \eqref{eq:gk-dim}.
  \item V\'an et al.\ \cite{Van2017} write the same term with $\ell^{2}$,
        where $\ell$ is a length, so that $\ell^{2}\equiv\kk$ dimensionally,
        and additionally report the dimensionless group
        $b=\ell^{2}/(\tau\alpha)$.
  \item Feh\'er and Kov\'acs \cite{FeherKovacs2021} tabulate $\kk$ directly in
        \si{\metre\squared} and quote the ratio $\kk/(\alpha\tq)$.
\end{itemize}

\begin{resultbox}
$b$ of \cite{Van2017} and $\kk/(\alpha\tq)$ of \cite{FeherKovacs2021} denote the
same dimensionless quantity, called $\Bnum$ in \cref{sec:Bdef}. No symbol from
one source is identified with a similarly named symbol from another without the
explicit conversion above. Where a source tabulates $\ell$ rather than $\kk$,
the value actually used in \cref{sec:calibrations} is the one the authors
themselves report, and any inconsistency is flagged rather than silently
recomputed.
\end{resultbox}

\subsection{Elimination of \texorpdfstring{$T$}{T}: single equation for the flux}

\begin{derivationx}[Flux equation]
\label{der:fluxeq}
Differentiate \eqref{eq:gk-dim} with respect to $x$:
\begin{equation}
  \tq\,\frac{\partial^{2} q}{\partial t\,\partial x}
  + \frac{\partial q}{\partial x}
  + \lambda\,\frac{\partial^{2} T}{\partial x^{2}}
  - \kk\,\frac{\partial^{3} q}{\partial x^{3}} = 0 .
  \label{eq:gk-dx}
\end{equation}
From the energy balance \eqref{eq:energy-dim},
\begin{equation}
  \frac{\partial T}{\partial t} = -\frac{1}{\rho c}\,
    \frac{\partial q}{\partial x}
  \qquad\Longrightarrow\qquad
  \frac{\partial^{2} T}{\partial x^{2}}
    = -\frac{1}{\rho c}\,\frac{\partial}{\partial t}
      \!\left(\frac{\partial q}{\partial x}\right)^{\!\!*}
  \label{eq:Txx}
\end{equation}
where the starred step uses the interchange of $\partial_{xx}$ and
$\partial_{t}$ applied to \eqref{eq:energy-dim} integrated once in $x$; more
directly, differentiating \eqref{eq:energy-dim} with respect to $x$ gives
$\rho c\,\partial_{t}\partial_{x}T = -\partial_{xx}q$, so
$\partial_{x}T$ satisfies
$\partial_{t}(\partial_{x}T) = -\partial_{xx}q/(\rho c)$.
Substituting $\lambda\,\partial_{xx}T$ from this relation into
\eqref{eq:gk-dx} is most cleanly done in the Laplace domain, which is why
\cref{sec:laplace} performs the elimination there. Carrying it out in the time
domain for zero initial data gives the equivalent single equation
\begin{equation}
  \boxed{\ \tq\,\frac{\partial^{2} q}{\partial t^{2}}
        + \frac{\partial q}{\partial t}
        = \alpha\,\frac{\partial^{2} q}{\partial x^{2}}
        + \kk\,\frac{\partial^{3} q}{\partial t\,\partial x^{2}}\ }
  \label{eq:gk-flux-dim}
\end{equation}
with $\alpha=\lambda/(\rho c)$.
\end{derivationx}

\begin{checkbox}
Dimensions of \eqref{eq:gk-flux-dim}, target
\si{\watt\per\metre\squared\per\second}:
$[\tq\partial_{tt}q]=\si{\second}\cdot\si{\watt\per\metre\squared\per\second\squared}
 =\si{\watt\per\metre\squared\per\second}$;
$[\partial_{t}q]=\si{\watt\per\metre\squared\per\second}$;
$[\alpha\partial_{xx}q]=\si{\metre\squared\per\second}\cdot
 \si{\watt\per\metre\squared\per\metre\squared}
 =\si{\watt\per\metre\squared\per\second}$;
$[\kk\partial_{t}\partial_{xx}q]=\si{\metre\squared}\cdot
 \si{\watt\per\metre\squared\per\second\per\metre\squared}
 =\si{\watt\per\metre\squared\per\second}$.
All four terms agree. Consistent.
\end{checkbox}

\begin{resultbox}
The same elimination performed on $T$ instead of $q$ produces an equation of
identical form in $T$, namely
$\tq\partial_{tt}T+\partial_{t}T=\alpha\partial_{xx}T+\kk\partial_{t}\partial_{xx}T$
\cite{Kovacs2018}. The two representations are \emph{not} interchangeable once
boundary data are imposed: see \cref{sec:representation}.
\end{resultbox}

\subsection{Symbol table}

\begin{table}[H]
\centering
\caption{Symbols, meanings, SI units and source.}
\label{tab:symbols}
\small
\begin{tabularx}{\textwidth}{@{}llX@{}}
\toprule
\textbf{Symbol} & \textbf{SI unit} & \textbf{Meaning and source} \\
\midrule
$T$        & \si{\kelvin} & temperature field, \cref{def:domain} \\
$q$        & \si{\watt\per\metre\squared} & heat flux ($x$-component) \\
$x$        & \si{\metre}  & position, $0\le x\le L$ \\
$t$        & \si{\second} & time \\
$L$        & \si{\metre}  & slab thickness \\
$t_{p}$    & \si{\second} & pulse duration; $t_{p}=\SI{0.01}{\second}$ in the
                            anchor experiments \cite{Kovacs2018} \\
$\rho$     & \si{\kilogram\per\metre\cubed} & mass density \\
$c$        & \si{\joule\per\kilogram\per\kelvin} & specific heat \\
$\rho c$   & \si{\joule\per\metre\cubed\per\kelvin} & volumetric heat capacity \\
$\lambda$  & \si{\watt\per\metre\per\kelvin} & thermal conductivity \\
$\alpha$   & \si{\metre\squared\per\second} & thermal diffusivity,
                                              $\alpha=\lambda/(\rho c)$ \\
$\tq$      & \si{\second} & flux relaxation time, \eqref{eq:gk-dim} \\
$\kk$      & \si{\metre\squared} & nonlocal coefficient, \eqref{eq:gk-dim} \\
$\Bnum$    & --- & resonance number $\kk/(\alpha\tq)$, \cref{sec:Bdef} \\
$q_{\max}$ & \si{\watt\per\metre\squared} & pulse amplitude \\
$\bar q_{0}$ & \si{\watt\per\metre\squared} & time-averaged pulse flux \\
$T_{0}$    & \si{\kelvin} & initial (ambient) temperature \\
$T_{\mathrm{end}}$ & \si{\kelvin} & adiabatic final temperature \\
$s$        & \si{\per\second} & Laplace variable conjugate to $t$ \\
$m(s)$     & \si{\per\metre}  & propagation factor, \eqref{eq:m2} \\
$\sigma$   & --- & noise standard deviation (dimensionless $T$) \\
\bottomrule
\end{tabularx}
\end{table}

%=======================================================================
\section{Initial and boundary conditions}
\label{sec:icbc}
%=======================================================================

\begin{definitionx}[Excitation]
\label{def:pulse}
The irradiated face receives a pulse of finite duration $t_{p}$ and cosine
shape,
\begin{equation}
  q(0,t) =
  \begin{cases}
    q_{\max}\left[1-\cos\!\left(\dfrac{2\pi t}{t_{p}}\right)\right],
      & 0 < t \leq t_{p}, \\[2ex]
    0, & t > t_{p} .
  \end{cases}
  \label{eq:pulse-dim}
\end{equation}
\end{definitionx}

\source{This is exactly the boundary function of the anchor
\cite{Kovacs2018}, which states $t_{p}=\SI{0.01}{\second}$ for the underlying
experiments \cite{Both2016,Van2017}. It is not invented here and is not
modified.}

\begin{definitionx}[Rear face and initial state]
\label{def:bcic}
\begin{align}
  q(L,t) &= 0 \qquad \text{for all } t \quad
            \text{(adiabatic rear face)} ,
  \label{eq:bc-rear-dim}\\
  T(x,0) &= T_{0} ,
  \label{eq:ic-T}\\
  q(x,0) &= 0 .
  \label{eq:ic-q}
\end{align}
Boundary cooling is neglected, as in the anchor \cite{Kovacs2018}.
\end{definitionx}

\subsection{The initial time-derivative is derived, not imposed}
\label{sec:ic-derived}

\begin{derivationx}[Consistency of the second initial condition]
\label{der:ic}
Equation \eqref{eq:gk-flux-dim} is second order in time, so a well-posed
problem needs both $q(x,0)$ and $\partial_{t}q(x,0)$. These cannot be
prescribed independently. Evaluating the constitutive law \eqref{eq:gk-dim} at
$t=0$ and solving for the time derivative,
\begin{equation}
  \tq\,\frac{\partial q}{\partial t}(x,0)
  = -\,q(x,0)
    - \lambda\,\frac{\partial T}{\partial x}(x,0)
    + \kk\,\frac{\partial^{2} q}{\partial x^{2}}(x,0) .
  \label{eq:ic-derived-eq}
\end{equation}
With the quiescent state \eqref{eq:ic-T}--\eqref{eq:ic-q}: $q(x,0)=0$ gives
$\partial_{xx}q(x,0)=0$, and $T(x,0)=T_{0}$ constant gives
$\partial_{x}T(x,0)=0$. Hence every term on the right vanishes and
\begin{equation}
  \frac{\partial q}{\partial t}(x,0) = 0 .
  \label{eq:ic-qdot}
\end{equation}
\end{derivationx}

\begin{resultbox}
$\partial_{t}q(x,0)=0$ is a \emph{consequence} of \eqref{eq:ic-T},
\eqref{eq:ic-q} and the constitutive law, not an independent assumption.
Imposing any other value would violate \eqref{eq:gk-dim} at $t=0$. The program
of \cref{sec:program} must not treat it as a free input.
\end{resultbox}

\subsection{Choice of representation}
\label{sec:representation}

\begin{assumptionx}[Flux representation]
\label{ass:qrep}
The calculations use the flux representation \eqref{eq:gk-flux-dim}, recovering
$T$ afterwards by integrating \eqref{eq:energy-dim} in time.
\end{assumptionx}

\textbf{Reason.} The experiment prescribes the \emph{flux} at $x=0$, equation
\eqref{eq:pulse-dim}. In the temperature representation this boundary datum is
not expressible in $T$ alone without introducing time derivatives of the
boundary temperature. The temperature representation is convenient only for
Dirichlet data in $T$. The anchor makes the same choice for the same reason
\cite{Kovacs2018,KovacsVan2015}.

%=======================================================================
\section{Nondimensionalisation}
\label{sec:nondim}
%=======================================================================

\subsection{Characteristic scales and substitutions}

\begin{definitionx}[Scaling]
\label{def:scaling}
\begin{align}
  \hat t &= \frac{\alpha t}{L^{2}}, &
  \hat x &= \frac{x}{L}, &
  \hat T &= \frac{T-T_{0}}{T_{\mathrm{end}}-T_{0}}, &
  \hat q &= \frac{q}{\bar q_{0}} .
  \label{eq:scaling}
\end{align}
\end{definitionx}

\source{Exactly the scaling of the anchor \cite{Kovacs2018}. Reproduced here in
full so every coefficient below can be checked independently.}

\begin{calculationx}[The averaged flux and the end temperature]
\label{calc:qbar}
The time-averaged pulse flux is
\begin{align}
  \bar q_{0}
   &= \frac{1}{t_{p}}\int_{0}^{t_{p}} q(0,t)\,\mathrm{d}t
    = \frac{q_{\max}}{t_{p}}\int_{0}^{t_{p}}
      \left[1-\cos\!\left(\frac{2\pi t}{t_{p}}\right)\right]\mathrm{d}t
      \nonumber\\
   &= \frac{q_{\max}}{t_{p}}
      \left[t - \frac{t_{p}}{2\pi}\sin\!\left(\frac{2\pi t}{t_{p}}\right)
      \right]_{0}^{t_{p}}
    = \frac{q_{\max}}{t_{p}}
      \left[t_{p} - \frac{t_{p}}{2\pi}\sin(2\pi)\right]
    = q_{\max} ,
  \label{eq:qbar}
\end{align}
since $\sin 2\pi = 0$. The energy deposited per unit area is
$\bar q_{0}t_{p}$; spread uniformly over the insulated slab of heat capacity
$\rho c L$ per unit area it raises the temperature by
$\bar q_{0}t_{p}/(\rho c L)$, so
\begin{equation}
  T_{\mathrm{end}} = T_{0} + \frac{\bar q_{0}\,t_{p}}{\rho c\,L} .
  \label{eq:Tend}
\end{equation}
\end{calculationx}

\begin{checkbox}
$[\bar q_{0}t_{p}/(\rho c L)]
 = \si{\watt\per\metre\squared}\cdot\si{\second}\cdot
   \si{\metre\cubed\kelvin\per\joule}\cdot\si{\per\metre}
 = \si{\joule\per\metre\squared}\cdot\si{\metre\squared\kelvin\per\joule}
 = \si{\kelvin}$. Correct. Consequently $\hat T\to1$ as $\hat t\to\infty$,
which is the normalisation used for all reported temperature histories.
Symbolic verification of \eqref{eq:qbar}: check~7d of
\texttt{check\_symbolic.py}, result $\bar q_{0}/q_{\max}=1$ exactly.
\end{checkbox}

\begin{definitionx}[Dimensionless parameters]
\label{def:ndparams}
\begin{align}
  \hat\tau_{\Delta} &= \frac{\alpha t_{p}}{L^{2}} , &
  \hat\tau_{q}      &= \frac{\alpha \tq}{L^{2}} , &
  \hat\kappa^{2}    &= \frac{\kk}{L^{2}} .
  \label{eq:ndparams}
\end{align}
\end{definitionx}

\begin{checkbox}
Each is dimensionless:
$[\alpha t_{p}/L^{2}]=\si{\metre\squared\per\second}\cdot\si{\second}\cdot
\si{\per\metre\squared}=1$;
$[\alpha\tq/L^{2}]=1$ identically;
$[\kk/L^{2}]=\si{\metre\squared}/\si{\metre\squared}=1$.
\end{checkbox}

\subsection{Line-by-line transformation of the field equations}

\begin{derivationx}[Energy balance in dimensionless variables]
\label{der:nd-energy}
Write $T = T_{0}+(T_{\mathrm{end}}-T_{0})\hat T$ and $q=\bar q_{0}\hat q$, with
$t=L^{2}\hat t/\alpha$ and $x=L\hat x$. Then
\begin{equation}
  \frac{\partial T}{\partial t}
   = (T_{\mathrm{end}}-T_{0})\,\frac{\partial \hat T}{\partial \hat t}\,
     \frac{\mathrm{d}\hat t}{\mathrm{d}t}
   = (T_{\mathrm{end}}-T_{0})\,\frac{\alpha}{L^{2}}\,
     \frac{\partial \hat T}{\partial \hat t} ,
  \qquad
  \frac{\partial q}{\partial x}
   = \frac{\bar q_{0}}{L}\,\frac{\partial \hat q}{\partial \hat x} .
  \label{eq:nd-chain}
\end{equation}
Substituting into \eqref{eq:energy-dim}:
\begin{equation}
  \rho c\,(T_{\mathrm{end}}-T_{0})\,\frac{\alpha}{L^{2}}\,
    \frac{\partial\hat T}{\partial\hat t}
  + \frac{\bar q_{0}}{L}\,\frac{\partial\hat q}{\partial\hat x} = 0 .
  \label{eq:nd-energy-raw}
\end{equation}
Insert $T_{\mathrm{end}}-T_{0}=\bar q_{0}t_{p}/(\rho c L)$ from
\eqref{eq:Tend} into the first coefficient:
\begin{equation}
  \rho c \cdot \frac{\bar q_{0}t_{p}}{\rho c L}\cdot\frac{\alpha}{L^{2}}
  = \frac{\bar q_{0}\,\alpha t_{p}}{L^{3}}
  = \frac{\bar q_{0}}{L}\cdot\frac{\alpha t_{p}}{L^{2}}
  = \frac{\bar q_{0}}{L}\,\hat\tau_{\Delta} .
  \label{eq:nd-coeff}
\end{equation}
So \eqref{eq:nd-energy-raw} becomes
$(\bar q_{0}/L)\left[\hat\tau_{\Delta}\partial_{\hat t}\hat T
 + \partial_{\hat x}\hat q\right]=0$, and dividing by
$\bar q_{0}/L\neq0$ gives, dropping hats,
\begin{equation}
  \boxed{\ \tau_{\Delta}\,\frac{\partial T}{\partial t}
        + \frac{\partial q}{\partial x} = 0\ }
  \label{eq:nd-energy}
\end{equation}
\end{derivationx}

\begin{derivationx}[Constitutive law in dimensionless variables]
\label{der:nd-gk}
Each term of \eqref{eq:gk-dim} transforms as
\begin{align}
  \tq\,\frac{\partial q}{\partial t}
   &= \tq\,\bar q_{0}\,\frac{\alpha}{L^{2}}\,
      \frac{\partial\hat q}{\partial\hat t}
    = \bar q_{0}\,\hat\tau_{q}\,\frac{\partial\hat q}{\partial\hat t} ,
  \label{eq:nd-t1}\\
  q &= \bar q_{0}\,\hat q ,
  \label{eq:nd-t2}\\
  \lambda\,\frac{\partial T}{\partial x}
   &= \lambda\,(T_{\mathrm{end}}-T_{0})\,\frac{1}{L}\,
      \frac{\partial\hat T}{\partial\hat x}
    = \lambda\,\frac{\bar q_{0}t_{p}}{\rho c L}\cdot\frac{1}{L}\,
      \frac{\partial\hat T}{\partial\hat x}
    = \bar q_{0}\,\frac{\alpha t_{p}}{L^{2}}\,
      \frac{\partial\hat T}{\partial\hat x}
    = \bar q_{0}\,\hat\tau_{\Delta}\,
      \frac{\partial\hat T}{\partial\hat x} ,
  \label{eq:nd-t3}\\
  \kk\,\frac{\partial^{2}q}{\partial x^{2}}
   &= \kk\,\frac{\bar q_{0}}{L^{2}}\,
      \frac{\partial^{2}\hat q}{\partial\hat x^{2}}
    = \bar q_{0}\,\hat\kappa^{2}\,
      \frac{\partial^{2}\hat q}{\partial\hat x^{2}} .
  \label{eq:nd-t4}
\end{align}
In \eqref{eq:nd-t3} the step $\lambda/(\rho c)=\alpha$ was used. Every term
carries the common factor $\bar q_{0}$; dividing it out and dropping hats,
\begin{equation}
  \boxed{\ \tq\,\frac{\partial q}{\partial t} + q
        + \tau_{\Delta}\,\frac{\partial T}{\partial x}
        - \kk\,\frac{\partial^{2} q}{\partial x^{2}} = 0\ }
  \label{eq:nd-gk}
\end{equation}
\end{derivationx}

\begin{checkbox}
Note that the coefficient of $\partial_{x}T$ in \eqref{eq:nd-gk} is
$\tau_{\Delta}$, not $1$ and not $\alpha$. This follows from
\eqref{eq:nd-t3} and depends on the choice $\hat T=(T-T_{0})/
(T_{\mathrm{end}}-T_{0})$ with $T_{\mathrm{end}}$ given by \eqref{eq:Tend}. A
different temperature normalisation would change this coefficient. The
appearance of the same $\tau_{\Delta}$ in \eqref{eq:nd-energy} and
\eqref{eq:nd-gk} is therefore not a coincidence but a consequence of the
scaling; the program must use both.
\end{checkbox}

\begin{derivationx}[Boundary and initial conditions, dimensionless]
\label{der:nd-bcic}
With $\hat q=q/\bar q_{0}$ and $\bar q_{0}=q_{\max}$ from \eqref{eq:qbar}, and
$t/t_{p} = \hat t/\hat\tau_{\Delta}$, the pulse \eqref{eq:pulse-dim} becomes
\begin{equation}
  q(0,t) =
  \begin{cases}
    1-\cos\!\left(\dfrac{2\pi t}{\tau_{\Delta}}\right),
      & 0 < t \leq \tau_{\Delta}, \\[2ex]
    0, & t > \tau_{\Delta},
  \end{cases}
  \qquad
  q(1,t) = 0 ,
  \label{eq:nd-bc}
\end{equation}
and the initial conditions are
\begin{equation}
  T(x,0)=0, \qquad q(x,0)=0, \qquad
  \frac{\partial q}{\partial t}(x,0)=0 ,
  \label{eq:nd-ic}
\end{equation}
the last by \cref{der:ic}. The domain is $0\le x\le 1$.
\end{derivationx}

\begin{derivationx}[Dimensionless flux equation]
\label{der:nd-flux}
Applying the same substitutions to \eqref{eq:gk-flux-dim}, or eliminating $T$
between \eqref{eq:nd-energy} and \eqref{eq:nd-gk},
\begin{equation}
  \boxed{\ \tq\,\frac{\partial^{2}q}{\partial t^{2}}
        + \frac{\partial q}{\partial t}
        = \frac{\partial^{2}q}{\partial x^{2}}
        + \kk\,\frac{\partial^{3}q}{\partial t\,\partial x^{2}}\ }
  \label{eq:nd-flux}
\end{equation}
The coefficient of $\partial_{xx}q$ is unity: the dimensionless diffusivity is
$\hat\alpha=\alpha t_{p}/L^{2}\big/\hat\tau_{\Delta}=1$ by construction.
\end{derivationx}

\begin{resultbox}
\textbf{The scaling sets $\hat\alpha\equiv1$.} The diffusivity is absorbed into
the time variable and cannot be read from the shape of a dimensionless curve;
it is recovered only through the physical time axis. This is essential when
converting between physical and dimensionless parameters in the program, and
was the source of a scaling error in earlier work that mixed
$\hat\alpha\equiv1$ with the physical $\alpha$ in the same expression.
\end{resultbox}

\begin{table}[H]
\centering
\caption{Dimensional to dimensionless mapping. $\alpha$ is eliminated from the
         dimensionless equations and reappears only through $\hat t$.}
\label{tab:mapping}
\small
\begin{tabular}{@{}llll@{}}
\toprule
\textbf{Dimensional} & \textbf{Characteristic scale} &
\textbf{Dimensionless} & \textbf{Equation} \\
\midrule
$t$        & $L^{2}/\alpha$                  & $\hat t=\alpha t/L^{2}$
           & \eqref{eq:scaling} \\
$x$        & $L$                             & $\hat x=x/L$
           & \eqref{eq:scaling} \\
$T-T_{0}$  & $T_{\mathrm{end}}-T_{0}$        & $\hat T$
           & \eqref{eq:scaling},\eqref{eq:Tend} \\
$q$        & $\bar q_{0}=q_{\max}$           & $\hat q$
           & \eqref{eq:scaling},\eqref{eq:qbar} \\
$t_{p}$    & $L^{2}/\alpha$                  &
  $\hat\tau_{\Delta}=\alpha t_{p}/L^{2}$     & \eqref{eq:ndparams} \\
$\tq$      & $L^{2}/\alpha$                  &
  $\hat\tau_{q}=\alpha \tq/L^{2}$            & \eqref{eq:ndparams} \\
$\kk$      & $L^{2}$                         &
  $\hat\kappa^{2}=\kk/L^{2}$                 & \eqref{eq:ndparams} \\
$\alpha$   & itself                          & $\hat\alpha\equiv1$
           & \cref{der:nd-flux} \\
$\Bnum$    & ---                             & invariant
           & \cref{sec:Bdef} \\
\bottomrule
\end{tabular}
\end{table}

%=======================================================================
\section{Laplace-domain formulation and the propagation factor}
\label{sec:laplace}
%=======================================================================

\begin{definitionx}[Laplace transform]
\label{def:laplace}
For a function $f(x,t)$ with $f(x,0^{-})=0$,
\begin{equation}
  \Lap f(x,s) = \int_{0}^{\infty} f(x,t)\,\mathrm{e}^{-st}\,\mathrm{d}t ,
  \qquad
  \mathcal{L}\!\left[\partial_{t}f\right] = s\Lap f - f(x,0) .
  \label{eq:laplace-def}
\end{equation}
A prime denotes $\partial/\partial x$ throughout this section.
\end{definitionx}

\begin{derivationx}[Transformed field equations]
\label{der:lap-fields}
Apply \eqref{eq:laplace-def} to the dimensionless energy balance
\eqref{eq:nd-energy} using $T(x,0)=0$:
\begin{equation}
  \tau_{\Delta}\left[s\Lap T - T(x,0)\right] + \Lap q^{\,\prime} = 0
  \quad\Longrightarrow\quad
  \tau_{\Delta}\,s\,\Lap T + \Lap q^{\,\prime} = 0 .
  \label{eq:lap-energy}
\end{equation}
Apply it to the constitutive law \eqref{eq:nd-gk} using $q(x,0)=0$:
\begin{equation}
  \tq\left[s\Lap q - q(x,0)\right] + \Lap q
    + \tau_{\Delta}\Lap T^{\,\prime} - \kk\,\Lap q^{\,\prime\prime} = 0
  \quad\Longrightarrow\quad
  \tq s\Lap q + \Lap q
    + \tau_{\Delta}\Lap T^{\,\prime} - \kk\Lap q^{\,\prime\prime} = 0 .
  \label{eq:lap-gk}
\end{equation}
\end{derivationx}

\begin{derivationx}[Elimination of $\Lap T$ and the propagation factor]
\label{der:m2}
Differentiate \eqref{eq:lap-energy} with respect to $x$:
\begin{equation}
  \tau_{\Delta}\,s\,\Lap T^{\,\prime} + \Lap q^{\,\prime\prime} = 0
  \quad\Longrightarrow\quad
  \tau_{\Delta}\,\Lap T^{\,\prime}
    = -\,\frac{\Lap q^{\,\prime\prime}}{s} .
  \label{eq:lap-Tprime}
\end{equation}
Substitute \eqref{eq:lap-Tprime} into \eqref{eq:lap-gk}:
\begin{equation}
  \tq s \Lap q + \Lap q
  - \frac{\Lap q^{\,\prime\prime}}{s}
  - \kk\,\Lap q^{\,\prime\prime} = 0 .
  \label{eq:lap-sub}
\end{equation}
Group the $\Lap q$ terms and the $\Lap q^{\,\prime\prime}$ terms:
\begin{equation}
  (1+\tq s)\,\Lap q
  = \left(\frac{1}{s} + \kk\right)\Lap q^{\,\prime\prime} .
  \label{eq:lap-group}
\end{equation}
Multiply both sides by $s$ and divide by $(1+\kk s)$, which is nonzero for
$s$ in the right half plane since $\kk\ge0$:
\begin{equation}
  \Lap q^{\,\prime\prime}
   = \frac{s\,(1+\tq s)}{1+\kk s}\,\Lap q .
  \label{eq:helmholtz-nd}
\end{equation}
Restoring $\alpha$ explicitly, so that the relation is valid in both the
dimensional and the dimensionless reading (with $\alpha=1$ in the latter),
\begin{equation}
  \boxed{\ \Lap q^{\,\prime\prime} = m^{2}(s)\,\Lap q ,
  \qquad
  m^{2}(s) = \frac{s\,(1+\tq s)}{\alpha + \kk s}\ }
  \label{eq:m2}
\end{equation}
\end{derivationx}

\begin{checkbox}[Origin of every factor in \eqref{eq:m2}]
\begin{itemize}[nosep]
  \item the factor $s$ in the numerator: the single time derivative in the
        energy balance \eqref{eq:nd-energy}, entering through
        \eqref{eq:lap-Tprime};
  \item the factor $(1+\tq s)$: the terms $\tq\partial_{t}q + q$ of the
        constitutive law \eqref{eq:nd-gk}, i.e. relaxation plus the flux
        itself;
  \item the term $\alpha$ in the denominator: the conduction term
        $\lambda\partial_{x}T$, which becomes $\tau_{\Delta}\partial_{x}T$ in
        \eqref{eq:nd-gk} and contributes the $1/s$ of \eqref{eq:lap-group};
  \item the term $\kk s$ in the denominator: the nonlocal term
        $-\kk\partial_{xx}q$ of \eqref{eq:nd-gk}.
\end{itemize}
\textbf{Symbolic verification.} Check 1a of \texttt{check\_symbolic.py} performs
the elimination independently with \texttt{sympy.solve} and returns
$s(1+\tq s)/(1+\kk s)$, matching \eqref{eq:helmholtz-nd} exactly. Check 1b
confirms that \eqref{eq:m2} reduces to it at $\alpha=1$. \NUM
\end{checkbox}

\begin{checkbox}[Dimensions of \eqref{eq:m2}]
$[\tq s] = \si{\second}\cdot\si{\per\second}=1$, so $(1+\tq s)$ is
dimensionally admissible. $[\kk s]=\si{\metre\squared}\cdot\si{\per\second}
=\si{\metre\squared\per\second}=[\alpha]$, so $(\alpha+\kk s)$ is admissible.
Hence
$[m^{2}] = \si{\per\second}\big/\si{\metre\squared\per\second}
 = \si{\per\metre\squared}$,
and $[m]=\si{\per\metre}$, correct for a quantity conjugate to $x$. Checks
2a--2c. \NUM
\end{checkbox}

\begin{definitionx}[Branch of $m(s)$]
\label{def:branch}
$m(s)=\sqrt{m^{2}(s)}$ is taken on the principal branch, $\operatorname{Re}m>0$.
Since $\alpha,\tq,\kk>0$, the zero of the numerator is at $s=-1/\tq$ and the
pole of the denominator at $s=-\alpha/\kk$, both on the negative real axis.
Therefore $m^{2}(s)$ is analytic and non-vanishing in the open right half
plane, and the principal branch is well defined there. This is the branch
required by the inversion contour of \cref{sec:numerics}.
\end{definitionx}

\begin{derivationx}[Spatial solution and the observable]
\label{der:That}
The general solution of \eqref{eq:m2} is
$\Lap q = A\cosh(mx)+C\sinh(mx)$. Imposing $\Lap q(0,s)=\Lap Q_{0}(s)$ and
$\Lap q(1,s)=0$ and solving the two-by-two system gives the compact form
\begin{equation}
  \Lap q(x,s)
   = \Lap Q_{0}(s)\,\frac{\sinh\!\left[m(1-x)\right]}{\sinh m} .
  \label{eq:qhat}
\end{equation}
Recover the temperature from \eqref{eq:lap-energy},
$\Lap T = -\Lap q^{\,\prime}/(\tau_{\Delta}s)$. Differentiating
\eqref{eq:qhat},
$\Lap q^{\,\prime} = -m\Lap Q_{0}\cosh\!\left[m(1-x)\right]/\sinh m$, so
\begin{equation}
  \boxed{\ \Lap T(x,s)
   = \frac{m\,\Lap Q_{0}(s)\,\cosh\!\left[m(1-x)\right]}
          {\tau_{\Delta}\,s\,\sinh m}\ }
  \label{eq:That}
\end{equation}
and at the observed rear face $x=1$, using $\cosh 0=1$,
\begin{equation}
  \boxed{\ \Lap T(1,s)
   = \frac{m\,\Lap Q_{0}(s)}{\tau_{\Delta}\,s\,\sinh m}\ }
  \label{eq:That-rear}
\end{equation}
\end{derivationx}

\begin{checkbox}
Checks 5c--5d of \texttt{check\_symbolic.py} verify by substitution that
\eqref{eq:qhat} satisfies \eqref{eq:m2} and both boundary conditions
$\Lap q(0)=\Lap Q_{0}$, $\Lap q(1)=0$, and that
$-\Lap q^{\,\prime}/(\tau_{\Delta}s)$ equals \eqref{eq:That}. \NUM
\end{checkbox}

\begin{derivationx}[Transform of the pulse]
\label{der:Q0}
With $\omega = 2\pi/\tau_{\Delta}$, direct integration of the finite cosine
pulse \eqref{eq:nd-bc} gives
\begin{align}
  \Lap Q_{0}(s)
   &= \int_{0}^{\tau_{\Delta}}
      \left[1-\cos\!\left(\frac{2\pi u}{\tau_{\Delta}}\right)\right]
      \mathrm{e}^{-su}\,\mathrm{d}u
  \nonumber\\
   &= \frac{\omega^{2}\left(1-\mathrm{e}^{-s\tau_{\Delta}}\right)}
           {s\left(s^{2}+\omega^{2}\right)} .
  \label{eq:Q0}
\end{align}
\end{derivationx}

\begin{checkbox}[Removable singularities --- numerically critical]
\label{check:removable}
The factor $s^{2}+\omega^{2}$ vanishes at $s=\pm\mathrm{i}\omega$. At those
points
\begin{equation}
  1-\mathrm{e}^{-s\tau_{\Delta}}\Big|_{s=\pm\mathrm{i}\omega}
   = 1-\mathrm{e}^{\mp 2\pi\mathrm{i}} = 1-1 = 0 ,
  \label{eq:removable}
\end{equation}
so the numerator vanishes simultaneously and the apparent poles are
\emph{removable}: $\Lap Q_{0}$ is analytic on the imaginary axis.
Checks 7a and 7b confirm \eqref{eq:Q0} and \eqref{eq:removable} symbolically.
\NUM

\textbf{Consequence for the program.} Any numerical scheme that splits
\eqref{eq:Q0} into $\omega^{2}/[s(s^{2}+\omega^{2})]$ times a shift operator
reinstates genuine poles at $s=\pm\mathrm{i}\omega$ and will fail. This is
quantified in \cref{sec:numerics} and is a mandatory requirement in the
program specification of \cref{sec:program}.
\end{checkbox}

\begin{resultbox}
The parameters $\tq$ and $\kk$ enter the observable \eqref{eq:That-rear}
\emph{only} through $m(s)$. The pulse transform $\Lap Q_{0}$ depends on neither.
This factorisation is what makes the analysis of \cref{sec:structural} exact and
independent of the pulse shape.
\end{resultbox}

%=======================================================================
\section{Definition and interpretation of \texorpdfstring{$\Bnum$}{B}}
\label{sec:Bdef}
%=======================================================================

\begin{derivationx}[Emergence of the ratio]
\label{der:Bemerge}
The denominator of \eqref{eq:m2} is $\alpha+\kk s$. Factor out $\alpha$:
\begin{equation}
  \alpha + \kk s
   = \alpha\left(1 + \frac{\kk}{\alpha}s\right)
   = \alpha\left(1 + \frac{\kk}{\alpha\tq}\,\tq s\right) .
  \label{eq:denom-factor}
\end{equation}
The numerator is $s(1+\tq s)$. Both now contain the same grouping $\tq s$, and
the only surviving free coefficient is the dimensionless combination
\begin{equation}
  \boxed{\ \Bnum = \frac{\kk}{\alpha\,\tq}\ }
  \label{eq:Bdef}
\end{equation}
so that
\begin{equation}
  m^{2}(s) = \frac{s\,(1+\tq s)}{\alpha\,(1+\Bnum\,\tq s)} .
  \label{eq:m2-B}
\end{equation}
\end{derivationx}

\begin{resultbox}[Mathematical meaning of $\Bnum$]
Equation \eqref{eq:m2-B} shows exactly what $\Bnum$ is: the ratio of the pole
location to the zero location of $m^{2}$ on the negative real axis. The
numerator vanishes at $\tq s=-1$; the denominator vanishes at
$\Bnum\tq s=-1$. $\Bnum$ measures how far the nonlocal contribution shifts the
pole relative to the relaxation zero. When $\Bnum=1$ the two coincide and
cancel; this is the content of \cref{sec:resonance}.
\end{resultbox}

\begin{checkbox}[Dimensions]
$[\Bnum]=\dfrac{[\kk]}{[\alpha][\tq]}
 =\dfrac{\si{\metre\squared}}
        {\si{\metre\squared\per\second}\cdot\si{\second}}
 =\dfrac{\si{\metre\squared}}{\si{\metre\squared}}=1$.
Dimensionless. Check 4a. \NUM
\end{checkbox}

\begin{calculationx}[Scale invariance: thickness does not change $\Bnum$]
\label{calc:Binv}
Evaluate $\Bnum$ in the dimensionless variables of \eqref{eq:ndparams},
remembering $\hat\alpha\equiv1$:
\begin{equation}
  \hat\Bnum
   = \frac{\hat\kappa^{2}}{\hat\alpha\,\hat\tau_{q}}
   = \frac{\kk/L^{2}}{1\cdot\left(\alpha\tq/L^{2}\right)}
   = \frac{\kk}{L^{2}}\cdot\frac{L^{2}}{\alpha\tq}
   = \frac{\kk}{\alpha\tq}
   = \Bnum .
  \label{eq:Binv}
\end{equation}
The thickness $L$ cancels identically.
\end{calculationx}

\begin{resultbox}
$\Bnum$ takes the same value in dimensional and dimensionless variables, and
all thicknesses of a given material share one value of $\Bnum$. Within this
1-D model, specimen thickness is \emph{not} a means of moving $\Bnum$.
Check 4d verifies \eqref{eq:Binv} symbolically. \NUM
\end{resultbox}

\subsection{Reparameterisation}
\label{sec:reparam-def}

\begin{definitionx}[Two parameter vectors]
\label{def:twovectors}
\begin{equation}
  \bm{\theta}=(\alpha,\ \tq,\ \kk) ,
  \qquad
  \bm{\phi}=(\alpha,\ \tq,\ \Bnum) .
  \label{eq:phi-def}
\end{equation}
\end{definitionx}

\begin{derivationx}[Forward and inverse maps]
\label{der:maps}
From \eqref{eq:Bdef}, the forward map $\bm{\theta}\mapsto\bm{\phi}$ is
$\alpha\mapsto\alpha$, $\tq\mapsto\tq$, $\Bnum=\kk/(\alpha\tq)$. Solving
\eqref{eq:Bdef} for $\kk$ gives the inverse map
\begin{equation}
  \boxed{\ \kk = \alpha\,\tq\,\Bnum\ }
  \label{eq:kappa-from-B}
\end{equation}
\end{derivationx}

\begin{calculationx}[Jacobian of the transformation]
\label{calc:jac-param}
With $\bm{\theta}=(\alpha,\tq,\kk)$ ordered as rows and
$\bm{\phi}=(\alpha,\tq,\Bnum)$ as columns, using \eqref{eq:kappa-from-B},
\begin{equation}
  \frac{\partial\bm{\theta}}{\partial\bm{\phi}}
  =
  \begin{pmatrix}
    \dfrac{\partial\alpha}{\partial\alpha} &
    \dfrac{\partial\alpha}{\partial\tq} &
    \dfrac{\partial\alpha}{\partial\Bnum} \\[1.6ex]
    \dfrac{\partial\tq}{\partial\alpha} &
    \dfrac{\partial\tq}{\partial\tq} &
    \dfrac{\partial\tq}{\partial\Bnum} \\[1.6ex]
    \dfrac{\partial\kk}{\partial\alpha} &
    \dfrac{\partial\kk}{\partial\tq} &
    \dfrac{\partial\kk}{\partial\Bnum}
  \end{pmatrix}
  =
  \begin{pmatrix}
    1 & 0 & 0 \\
    0 & 1 & 0 \\
    \tq\Bnum & \alpha\Bnum & \alpha\tq
  \end{pmatrix} ,
  \label{eq:param-jacobian}
\end{equation}
whose determinant, expanding along the first row, is
\begin{equation}
  \det\frac{\partial\bm{\theta}}{\partial\bm{\phi}}
   = 1\cdot\left(1\cdot\alpha\tq - 0\right) = \alpha\,\tq .
  \label{eq:param-det}
\end{equation}
\end{calculationx}

\begin{resultbox}
$\det(\partial\bm{\theta}/\partial\bm{\phi})=\alpha\tq\neq0$ for all
$\alpha,\tq>0$. The reparameterisation is therefore a smooth bijection on the
positive octant and creates or destroys no information anywhere, including at
$\Bnum=1$. Any degeneracy found in \cref{sec:structural} is a property of the
map from parameters to observations, not an artefact of the coordinates.
Checks 4b and 4c verify \eqref{eq:kappa-from-B} and \eqref{eq:param-det}
symbolically. \NUM

This change of variables is an elementary device used to organise the analysis.
It is not presented as a result.
\end{resultbox}

%=======================================================================
\section{Limiting cases}
\label{sec:limits}
%=======================================================================

All three limits are taken in the \emph{same} governing relation \eqref{eq:m2},
so that the four models compared in \cref{sec:bic} differ only by constraints on
$(\tq,\kk)$ and share one solver.

\begin{derivationx}[Fourier limit]
\label{der:fourier}
Set $\tq\to0$ and $\kk\to0$ in \eqref{eq:m2}:
\begin{equation}
  m^{2}(s)
   = \frac{s(1+\tq s)}{\alpha+\kk s}
   \;\xrightarrow[\ \kk\to0\ ]{\ \tq\to0\ }\;
   \frac{s(1+0)}{\alpha+0}
   = \frac{s}{\alpha} .
  \label{eq:lim-fourier}
\end{equation}
Inverting $\Lap q^{\,\prime\prime}=(s/\alpha)\Lap q$ corresponds to
$\partial_{t}T=\alpha\partial_{xx}T$, the diffusion equation.
\end{derivationx}

\textbf{Interpretation.} One parameter, $\alpha$. Infinite propagation speed,
purely diffusive rise. \textbf{Numerical check:} setting $\tq,\kk\to0$ in the
solver reproduces the analytical Fourier solution to $2.6\times10^{-6}$
\PREV.

\begin{derivationx}[Maxwell--Cattaneo--Vernotte limit]
\label{der:mcv}
Set $\kk\to0$ with $\tq>0$:
\begin{equation}
  m^{2}(s)
   \;\xrightarrow[\ \kk\to0\ ]{}\;
   \frac{s(1+\tq s)}{\alpha}
   = \frac{s}{\alpha} + \frac{\tq s^{2}}{\alpha} .
  \label{eq:lim-mcv}
\end{equation}
The $s^{2}$ term is the transform of $\partial_{tt}$, so this corresponds to
$\tq\partial_{tt}T+\partial_{t}T=\alpha\partial_{xx}T$, a damped wave equation.
\end{derivationx}

\textbf{Interpretation and wave speed.} As $s\to\infty$,
$m^{2}\sim\tq s^{2}/\alpha$, so $m\sim s\sqrt{\tq/\alpha}$ and the
characteristic $\mathrm{e}^{-mx}$ propagates a front at finite speed
\begin{equation}
  v = \sqrt{\alpha/\tq} .
  \label{eq:mcv-speed}
\end{equation}
Check 3d verifies $\lim_{s\to\infty}m^{2}/s^{2}=\tq/\alpha$ symbolically. \NUM
Two parameters, $(\alpha,\tq)$.

\begin{derivationx}[Nyíri limit]
\label{der:nyiri}
Set $\tq\to0$ with $\kk>0$:
\begin{equation}
  m^{2}(s)
   \;\xrightarrow[\ \tq\to0\ ]{}\;
   \frac{s}{\alpha+\kk s} .
  \label{eq:lim-nyiri}
\end{equation}
This retains the nonlocal term and discards relaxation.
\end{derivationx}

\source{The entropy-current formulation underlying this reduced form is due to
Nyíri \cite{Nyiri1991}; it is used here solely as the constraint $\tq=0$ on the
same transfer function, because it isolates the mechanism that the
Maxwell--Cattaneo--Vernotte limit omits. Two parameters, $(\alpha,\kk)$.}

\begin{derivationx}[Full Guyer--Krumhansl model]
\label{der:gkfull}
No constraint: $m^{2}(s)=s(1+\tq s)/(\alpha+\kk s)$, three parameters
$(\alpha,\tq,\kk)$.
\end{derivationx}

\begin{checkbox}
Dimensions of every limit: each of \eqref{eq:lim-fourier},
\eqref{eq:lim-mcv} and \eqref{eq:lim-nyiri} has dimension
$\si{\per\metre\squared}$, since $[s/\alpha]=\si{\per\second}\cdot
\si{\second\per\metre\squared}=\si{\per\metre\squared}$ and
$[\tq s^{2}/\alpha]=\si{\second}\cdot\si{\per\second\squared}\cdot
\si{\second\per\metre\squared}=\si{\per\metre\squared}$. Checks 3a--3c verify
all three limits symbolically. \NUM
\end{checkbox}

\begin{table}[H]
\centering
\caption{The four models as constraints on one transfer function
         \eqref{eq:m2}. $k$ is the number of free parameters, used in the
         BIC penalty of \cref{sec:bic}.}
\label{tab:models}
\small
\begin{tabular}{@{}llll@{}}
\toprule
\textbf{Model} & \textbf{Constraint} & $m^{2}(s)$ & $k$ \\
\midrule
Fourier & $\tq=0$, $\kk=0$ & $s/\alpha$ & 1 \\
MCV     & $\kk=0$          & $s(1+\tq s)/\alpha$ & 2 \\
Nyíri   & $\tq=0$          & $s/(\alpha+\kk s)$  & 2 \\
GK      & none             & $s(1+\tq s)/(\alpha+\kk s)$ & 3 \\
\bottomrule
\end{tabular}
\end{table}

MGT is excluded from this set: it requires a thermal-displacement variable that
cannot be written as a constraint on \eqref{eq:m2}, so including it would break
the shared-solver protocol.

%=======================================================================
\section{The Fourier-resonance condition}
\label{sec:resonance}
%=======================================================================

\begin{propositionx}[Cancellation at $\Bnum=1$]
\label{prop:cancel}
For $\alpha,\tq>0$, the condition $\Bnum=1$ is equivalent to
$\kk=\alpha\tq$, and under it $m^{2}(s)$ reduces identically to the Fourier
propagation factor.
\end{propositionx}

\begin{proof}
From \eqref{eq:Bdef}, $\Bnum=1 \iff \kk/(\alpha\tq)=1 \iff \kk=\alpha\tq$.
Substitute $\kk=\alpha\tq$ into the denominator of \eqref{eq:m2}, one step at a
time:
\begin{align}
  \alpha + \kk s
   &= \alpha + (\alpha\tq)\,s
  \label{eq:res-step1}\\
   &= \alpha\,(1 + \tq s) .
  \label{eq:res-step2}
\end{align}
Therefore
\begin{align}
  m^{2}(s)\Big|_{\Bnum=1}
   &= \frac{s\,(1+\tq s)}{\alpha\,(1+\tq s)}
  \label{eq:res-step3}\\
   &= \frac{s}{\alpha}\cdot\frac{(1+\tq s)}{(1+\tq s)}
  \label{eq:res-step4}\\
   &= \frac{s}{\alpha} ,
  \label{eq:res-step5}
\end{align}
the cancellation in \eqref{eq:res-step4} being legitimate for every $s$ because
$\tq>0$ implies $(1+\tq s)\not\equiv0$. Comparing \eqref{eq:res-step5} with the
Fourier limit \eqref{eq:lim-fourier} shows the two are identical.
\end{proof}

\begin{checkbox}
Check 5a verifies the factorisation \eqref{eq:res-step2}
($\alpha+\kk s$ with $\kk=\alpha\tq$ factors to $\alpha(1+\tq s)$), and check
5b verifies $m^{2}|_{\Bnum=1}-s/\alpha=0$ exactly. \NUM
\end{checkbox}

\begin{resultbox}
On the ray $\kk=\alpha\tq$ the Guyer--Krumhansl propagation factor is
\emph{exactly} the Fourier one, for every $s$, and the relaxation time has
disappeared from it entirely. The single factor $(1+\tq s)$ carried the whole
dependence of $m^{2}$ on $\tq$; when it cancels, $\tq$ leaves the problem.
\end{resultbox}

\source{\textbf{The Fourier-resonance condition is not claimed as a discovery
of this study.} It is established in the literature. Fülöp, Kovács and Ván
describe a hierarchical resonance at which the solutions of the coupled system
may be identical to those of the single Fourier equation
\cite{FulopKovacsVan2015}. The validation anchor states it explicitly and names
it, writing that the GK equation ``is able to recover the solution of Fourier
equation when $\kappa^{2}/\tau=\alpha$ holds, called Fourier resonance''
\cite{Kovacs2018}. Fülöp et al.\ describe the Fourier solutions as being
covered by those of the GK equation \cite{Fulop2018entropy}. Feh\'er and
Kov\'acs use the same ratio as the lower limit of their fitting domain
\cite{FeherKovacs2021}. The calculations of \cref{sec:structural} concern the
\emph{consequences} of this known condition for the inverse problem.}

%=======================================================================
\section{Exact structural identifiability}
\label{sec:structural}
%=======================================================================

\begin{definitionx}[Structural identifiability]
\label{def:struct}
Let $\mathcal{M}(\bm{\theta})$ be the map from parameters to the noise-free
observable $\{T(1,t):t>0\}$. A parameter or combination is structurally
identifiable if
$\mathcal{M}(\bm{\theta}_{1})=\mathcal{M}(\bm{\theta}_{2})$ for all $t$ implies
equality of that parameter in $\bm{\theta}_{1}$ and $\bm{\theta}_{2}$.
Structural non-identifiability is a property of the model and the observation
operator alone: no reduction of noise and no increase of sampling density can
remove it.
\end{definitionx}

\subsection{Complete dependence of the forward solution}

\begin{derivationx}[The response on the resonance ray]
\label{der:resp-ray}
By \eqref{eq:That}, the full dependence of the observable on $(\tq,\kk)$ is
through $m(s)$ alone; $\Lap Q_{0}$, $\tau_{\Delta}$, $s$ and $x$ do not involve
them. By \cref{prop:cancel}, on the ray
\begin{equation}
  m\big|_{\Bnum=1} = \sqrt{\frac{s}{\alpha}} ,
  \label{eq:m-ray}
\end{equation}
which contains no $\tq$. Substituting \eqref{eq:m-ray} into \eqref{eq:That},
\begin{equation}
  \Lap T(x,s)\Big|_{\Bnum=1}
   = \frac{\sqrt{s/\alpha}\;\Lap Q_{0}(s)\,
           \cosh\!\left[\sqrt{s/\alpha}\,(1-x)\right]}
          {\tau_{\Delta}\,s\,\sinh\!\sqrt{s/\alpha}} .
  \label{eq:That-ray}
\end{equation}
Every symbol on the right is independent of $\tq$. Hence
\begin{equation}
  \frac{\partial}{\partial \tq}\,
  \Lap T(x,s)\Big|_{\Bnum=1} \equiv 0
  \qquad \forall s,\ \forall x\in[0,1] .
  \label{eq:dTdtq-ray}
\end{equation}
The Laplace transform is injective on continuous functions, so the same holds
for $T(x,t)$ at every $t$.
\end{derivationx}

\begin{checkbox}
Check 5c confirms $\partial_{\tq}m|_{\text{ray}}=0$, and check 5d differentiates
the \emph{full} expression \eqref{eq:That-ray} with respect to $\tq$ at
arbitrary $x$ and returns exactly $0$. \NUM
\end{checkbox}

\begin{resultbox}[Non-injectivity]
Varying $\tq$ while holding $\kk=\alpha\tq$ traces the one-parameter family
\begin{equation}
  \mathcal{R}=\left\{(\alpha,\ \tq,\ \alpha\tq)\ :\ \tq>0\right\}
  \label{eq:family}
\end{equation}
along which the observable is constant. $\mathcal{M}$ is therefore not
injective, and by \cref{def:struct} the pair $(\tq,\kk)$ is structurally
non-identifiable on $\mathcal{R}$. This is not a coincidence of two particular
solutions but an entire continuum of indistinguishable parameter sets.
\end{resultbox}

\subsection{Sensitivities}

\begin{derivationx}[Derivatives of $m$ with respect to the parameters]
\label{der:dm}
Write $m=(N/D)^{1/2}$ with $N=s(1+\tq s)$ and $D=\alpha+\kk s$. Then
$\partial_{p}m = \tfrac{1}{2}m^{-1}\,\partial_{p}(N/D)$.

For $p=\tq$: $\partial_{\tq}N = s\cdot s = s^{2}$ and $\partial_{\tq}D=0$, so
$\partial_{\tq}(N/D)=s^{2}/D$ and
\begin{equation}
  \frac{\partial m}{\partial \tq}
   = \frac{1}{2m}\cdot\frac{s^{2}}{\alpha+\kk s}
   = \frac{s^{2}}{2\,m\,(\alpha+\kk s)} .
  \label{eq:dm-dtq}
\end{equation}

For $p=\kk$: $\partial_{\kk}N=0$ and $\partial_{\kk}D=s$, so
$\partial_{\kk}(N/D) = -N s/D^{2} = -m^{2}s/D$ and
\begin{equation}
  \frac{\partial m}{\partial \kk}
   = \frac{1}{2m}\cdot\left(-\frac{m^{2}s}{\alpha+\kk s}\right)
   = -\,\frac{m\,s}{2\,(\alpha+\kk s)} .
  \label{eq:dm-dk2}
\end{equation}
\end{derivationx}

\begin{checkbox}
Checks 5e and 5f verify \eqref{eq:dm-dtq} and \eqref{eq:dm-dk2} against direct
symbolic differentiation of $m$. \NUM
\end{checkbox}

\begin{derivationx}[Sensitivity ratio]
\label{der:ratio}
Because $\tq$ and $\kk$ enter $\Lap T$ only through $m$, the chain rule gives
$\partial_{p}\Lap T = (\partial_{m}\Lap T)(\partial_{p}m)$ for
$p\in\{\tq,\kk\}$. Wherever $\partial_{m}\Lap T\neq0$ the common factor cancels
in a ratio:
\begin{equation}
  \frac{\partial \Lap T/\partial \tq}
       {\partial \Lap T/\partial \kk}
  = \frac{\partial m/\partial \tq}{\partial m/\partial \kk} .
  \label{eq:ratio-chain}
\end{equation}
Dividing \eqref{eq:dm-dtq} by \eqref{eq:dm-dk2}, the factor
$2(\alpha+\kk s)$ cancels:
\begin{equation}
  \frac{\partial m/\partial \tq}{\partial m/\partial \kk}
   = \frac{s^{2}/m}{-\,m\,s}
   = -\,\frac{s}{m^{2}} .
  \label{eq:ratio-general}
\end{equation}
This holds for \emph{all} parameter values. Now evaluate on the ray, where
$m^{2}=s/\alpha$ by \cref{prop:cancel}:
\begin{equation}
  \boxed{\ \frac{\partial \Lap T/\partial \tq}
       {\partial \Lap T/\partial \kk}\Bigg|_{\Bnum=1}
   = -\,\frac{s}{s/\alpha}
   = -\,\alpha\ }
  \label{eq:ratio-resonance}
\end{equation}
\end{derivationx}

\begin{resultbox}
The ratio \eqref{eq:ratio-resonance} is a \emph{constant}: independent of $s$,
hence independent of $t$ after inversion; independent of $x$; and independent
of the pulse, since $\Lap Q_{0}$ cancelled in \eqref{eq:ratio-chain}. Therefore
in the time domain, pointwise,
\begin{equation}
  \frac{\partial T}{\partial \tq}
   = -\,\alpha\,\frac{\partial T}{\partial \kk} .
  \label{eq:antiparallel}
\end{equation}
The two sensitivity fields are exactly anti-parallel. Checks 5g and 5h verify
\eqref{eq:ratio-general} and \eqref{eq:ratio-resonance}. \NUM
\end{resultbox}

\subsection{Jacobian rank, Fisher singularity, correlation}

\begin{definitionx}[Sensitivity matrix]
\label{def:jacobian}
For observations $y_{i}=T(1,t_{i})$, $i=1,\dots,n$, and parameters
$\theta_{j}$,
\begin{equation}
  J_{ij} = \frac{\partial y_{i}}{\partial \theta_{j}} .
  \label{eq:J-def}
\end{equation}
\end{definitionx}

\begin{derivationx}[Rank deficiency of the $(\tq,\kk)$ block]
\label{der:rank}
Restrict to the two columns $j\in\{\tq,\kk\}$ and evaluate on the ray. By
\eqref{eq:antiparallel}, for every observation index $i$,
\begin{equation}
  J_{i,\tq} = -\alpha\, J_{i,\kk} ,
  \label{eq:col-dependence}
\end{equation}
that is, the first column is a constant multiple of the second:
\begin{equation}
  J_{:,\tq} = -\alpha\, J_{:,\kk} .
  \label{eq:col-vector}
\end{equation}
Two linearly dependent columns give
\begin{equation}
  \operatorname{rank}\left[J_{:,\tq}\ \ J_{:,\kk}\right] = 1 ,
  \label{eq:rank1}
\end{equation}
and the null direction follows from
$J_{:,\tq}\cdot 1 + J_{:,\kk}\cdot\alpha
 = -\alpha J_{:,\kk} + \alpha J_{:,\kk} = \bm{0}$, i.e.
\begin{equation}
  \left[J_{:,\tq}\ \ J_{:,\kk}\right]
  \begin{pmatrix}1\\ \alpha\end{pmatrix} = \bm{0} .
  \label{eq:nullvec}
\end{equation}
\end{derivationx}

\begin{resultbox}[Geometric meaning of the null vector]
The direction $(1,\alpha)$ in the $(\tq,\kk)$ plane is precisely the tangent to
the ray $\kk=\alpha\tq$, since along it $\mathrm{d}\kk/\mathrm{d}\tq=\alpha$.
The flat direction of the objective function \emph{is} the resonance ray.
\end{resultbox}

\begin{derivationx}[Fisher information and its singularity]
\label{der:fisher}
Assume additive independent Gaussian noise of constant variance $\sigma^{2}$,
$y_{i}^{\text{obs}} = y_{i}(\bm{\theta}) + \varepsilon_{i}$,
$\varepsilon_{i}\sim\mathcal{N}(0,\sigma^{2})$. The log-likelihood is
\begin{equation}
  \ln \mathcal{L}
   = \text{const} - \frac{1}{2\sigma^{2}}
     \sum_{i=1}^{n}\left[y_{i}^{\text{obs}}-y_{i}(\bm{\theta})\right]^{2} ,
  \label{eq:loglik}
\end{equation}
and the Fisher information matrix is
\begin{equation}
  F_{jk}
   = \mathbb{E}\!\left[
     \frac{\partial \ln\mathcal{L}}{\partial\theta_{j}}
     \frac{\partial \ln\mathcal{L}}{\partial\theta_{k}}\right]
   = \frac{1}{\sigma^{2}}\sum_{i=1}^{n}
     \frac{\partial y_{i}}{\partial\theta_{j}}
     \frac{\partial y_{i}}{\partial\theta_{k}}
   \qquad\Longleftrightarrow\qquad
   F = \frac{J^{\!\top}J}{\sigma^{2}} .
  \label{eq:fisher-def}
\end{equation}
On the ray, substituting \eqref{eq:col-vector} into the two-by-two block,
\begin{equation}
  F^{(\tq,\kk)}
  = \frac{1}{\sigma^{2}}
    \begin{pmatrix}
      \alpha^{2}\,\|J_{:,\kk}\|^{2} & -\alpha\,\|J_{:,\kk}\|^{2} \\
      -\alpha\,\|J_{:,\kk}\|^{2}    & \|J_{:,\kk}\|^{2}
    \end{pmatrix}
  = \frac{\|J_{:,\kk}\|^{2}}{\sigma^{2}}
    \begin{pmatrix} \alpha^{2} & -\alpha \\ -\alpha & 1 \end{pmatrix} ,
  \label{eq:F-block}
\end{equation}
whose determinant is
\begin{equation}
  \det F^{(\tq,\kk)}
   = \left(\frac{\|J_{:,\kk}\|^{2}}{\sigma^{2}}\right)^{2}
     \left[\alpha^{2}\cdot 1 - (-\alpha)(-\alpha)\right]
   = \left(\frac{\|J_{:,\kk}\|^{2}}{\sigma^{2}}\right)^{2}
     \left[\alpha^{2}-\alpha^{2}\right]
   = 0 .
  \label{eq:detF}
\end{equation}
The implied parameter correlation is
\begin{equation}
  \operatorname{corr}(\tq,\kk)
   = \frac{F_{12}}{\sqrt{F_{11}F_{22}}}
   = \frac{-\alpha}{\sqrt{\alpha^{2}\cdot 1}}
   = -1 .
  \label{eq:corr}
\end{equation}
\end{derivationx}

\begin{propositionx}[Structural non-identifiability at $\Bnum=1$]
\label{prop:main}
For the system of \cref{sec:problem} with $\kk=\alpha\tq$:
\begin{enumerate}[nosep,label=(\roman*)]
  \item $T(x,t)$ is independent of $\tq$ at every $x\in[0,1]$ and every $t$,
        for any pulse;
  \item the family $\mathcal{R}$ of \eqref{eq:family} produces identical
        observations, so $(\tq,\kk)$ are not separately structurally
        identifiable;
  \item $\partial T/\partial\tq = -\alpha\,\partial T/\partial\kk$ exactly, the
        sensitivity matrix has rank one with null direction $(1,\alpha)$, the
        Fisher matrix is singular, and $\operatorname{corr}(\tq,\kk)=-1$.
\end{enumerate}
\end{propositionx}

\begin{proof}
(i) is \cref{der:resp-ray}, equation \eqref{eq:dTdtq-ray}. (ii) follows from (i)
by \cref{def:struct}. (iii) is \eqref{eq:antiparallel}, \eqref{eq:rank1},
\eqref{eq:nullvec}, \eqref{eq:detF} and \eqref{eq:corr}.
\end{proof}

\begin{checkbox}[Independent symbolic verification of the whole proposition]
Checks 6a--6d of \texttt{check\_symbolic.py} construct the Jacobian row
$\left[g\,\partial_{\tq}m,\ g\,\partial_{\kk}m\right]$ with $g=\partial_{m}T$
symbolic and nonzero, evaluate on $\kk=\alpha\tq$, and confirm: the entry ratio
is $-\alpha$; $(1,\alpha)$ annihilates the row; $\det(J^{\!\top}J)=0$; the
product of the two entries is negative, giving correlation $-1$.

\emph{Technical note recorded for reproducibility.} Sympy does not
automatically reduce $\sqrt{s^{2}\tq^{2}+2s\tq+1}$ to $(1+s\tq)$. The
polynomial factors exactly as $(1+s\tq)^{2}$ and both factors are positive for
$s,\tq>0$, so the reduction is valid; the check script applies it explicitly.
Numerical spot checks confirm the difference is zero to machine precision.
Without this reduction the checks return unsimplified expressions rather than
$0$; this is a limitation of the algebra system, not of the result. \NUM
\end{checkbox}

\begin{resultbox}[What is \emph{not} claimed]
\cref{prop:main} concerns the pair $(\tq,\kk)$ only. It does \emph{not} state
that GK parameters are unidentifiable in general. By \eqref{eq:That-ray} the
response on the ray is a Fourier response with diffusivity $\alpha$, so
$\alpha$ remains determined by the physical time axis, and the statement
$\Bnum=1$ is itself an identified fact. See \cref{sec:reparam}.
\end{resultbox}

%=======================================================================
\section{Numerical proof that the degeneracy is exact}
\label{sec:numdegen}
%=======================================================================

The purpose of this section is to exclude the possibility that
\cref{prop:main} is an artefact of algebraic manipulation, and to exclude
discretisation, interpolation, inversion and round-off as explanations.

\begin{calculationx}[Setup]
\label{calc:degen-setup}
\textbf{Parameter setup.} Three parameter sets on the ray, with $\tq$ spanning
two orders of magnitude and $\kk$ slaved to it:
\begin{equation}
  \tq \in \{3\times10^{-3},\ 3\times10^{-2},\ 3\times10^{-1}\},
  \qquad \kk = \alpha\tq \ \ (\Bnum=1), \qquad \alpha=1 .
  \label{eq:degen-params}
\end{equation}
\textbf{Control.} The identical three $\tq$ values with $\kk=1.05\,\alpha\tq$,
i.e. $\Bnum=1.05$, a $5\%$ departure from resonance.

\textbf{Observation.} Rear face $x=1$, on the grid
$t\in[0.005,1.0]$ with $400$ points, $\tau_{\Delta}=0.04$.

\textbf{Comparison metric.} The pointwise spread across the three curves,
maximised over the grid:
\begin{equation}
  S = \max_{t}\left[\max_{j}T_{j}(t) - \min_{j}T_{j}(t)\right] ,
  \qquad j=1,2,3 .
  \label{eq:spread}
\end{equation}
\textbf{Expected result.} If \cref{prop:main} holds, $S$ on the ray is limited
only by the numerical noise floor of the solver, whereas the control must show
a spread of order the physical difference between $\Bnum=1$ and $\Bnum=1.05$.
\end{calculationx}

\begin{resultbox}[Double precision, certified convolution solver]
\begin{align}
  S_{\text{ray}} &= 1.261\times10^{-11} , &&\NUM
  \label{eq:S-ray}\\
  S_{\text{control}} &= 2.595\times10^{-2} , &&\NUM
  \label{eq:S-ctl}\\
  S_{\text{control}}/S_{\text{ray}} &= 2.06\times10^{9} . &&\NUM
  \label{eq:S-ratio}
\end{align}
Script \texttt{validation/calc/check\_numerical.py}, checks 7a--7c.
\end{resultbox}

\begin{calculationx}[Independent arbitrary-precision confirmation]
\label{calc:degen-mp}
To exclude double-precision round-off entirely, the same test was repeated in
$50$-digit arithmetic using an independent implementation
(\texttt{mpmath}, convolution of the kernel with the pulse, Talbot inversion at
degree $30$), at the single probe $t=0.5$:
\begin{equation}
  S^{(50\text{-digit})}_{\text{ray}} = 7.733\times10^{-30} . \qquad\NUM
  \label{eq:S-mp}
\end{equation}
Script \texttt{validation/calc/check\_numerical.py}, check 7d.
\end{calculationx}

\begin{checkbox}
The spread falls from $1.3\times10^{-11}$ in double precision to
$7.7\times10^{-30}$ at $50$ digits, i.e. it tracks the working precision rather
than settling at a fixed nonzero value. That is the signature of an exactly
zero quantity computed approximately, not of a small but genuine difference.
The control at $\Bnum=1.05$ does not collapse and is nine orders of magnitude
larger, confirming that the effect is specific to the ray and is not a
consequence of the solver, the grid or the observation point.
\end{checkbox}

\begin{resultbox}
Ruled out as explanations: discretisation, interpolation, Laplace inversion,
finite precision and coding error. The degeneracy is exact, consistent with
\cref{prop:main}.
\end{resultbox}

%=======================================================================
\section{Sensitivity matrix and Fisher information}
\label{sec:fisher}
%=======================================================================

\begin{definitionx}[Observation model]
\label{def:obsmodel}
\begin{equation}
  y_{i}^{\text{obs}} = T(1,t_{i};\bm{\theta}) + \varepsilon_{i} ,
  \qquad
  \varepsilon_{i}\ \text{i.i.d.}\ \mathcal{N}(0,\sigma^{2}),
  \qquad i=1,\dots,n .
  \label{eq:obsmodel}
\end{equation}
The noise level is specified as a fraction $\eta$ of the temperature range:
\begin{equation}
  \sigma = \eta\left[\max_{i}T(1,t_{i}) - \min_{i}T(1,t_{i})\right] .
  \label{eq:sigma-def}
\end{equation}
\end{definitionx}

\begin{definitionx}[Log-parameterisation]
\label{def:logparam}
Estimation is performed in $\ln\theta_{j}$, which enforces positivity and makes
the three parameters comparably scaled. The chain rule gives the log-Jacobian
\begin{equation}
  \tilde J_{ij}
   = \frac{\partial y_{i}}{\partial \ln\theta_{j}}
   = \theta_{j}\,\frac{\partial y_{i}}{\partial \theta_{j}}
   = \theta_{j}\,J_{ij}
   \quad\text{(no sum)} ,
  \label{eq:logJ}
\end{equation}
so standard errors obtained from $\tilde F=\tilde J^{\!\top}\tilde J/\sigma^{2}$
are \emph{relative} standard errors, reported below as percentages.
\end{definitionx}

\begin{derivationx}[Analytic sensitivities used by the program]
\label{der:analytic-sens}
Finite differencing of the time-domain solver is not used anywhere in the
statistics, because the solver noise floor contaminates the Jacobian. Instead,
writing the rear-face kernel of \eqref{eq:factorisation} below as
$\Lap K = m/(\tau_{\Delta}s\sinh m)$, differentiate with respect to $m$ using
the quotient rule:
\begin{align}
  \frac{\partial \Lap K}{\partial m}
   &= \frac{1}{\tau_{\Delta}s}\cdot
      \frac{1\cdot\sinh m - m\cosh m}{\sinh^{2}m}
  \label{eq:dKdm-1}\\
   &= \frac{1}{\tau_{\Delta}s\,\sinh m}
      \left(1 - m\,\frac{\cosh m}{\sinh m}\right)
  \label{eq:dKdm-2}\\
   &= \frac{1}{\tau_{\Delta}s\,\sinh m}\left(1 - m\coth m\right) .
  \label{eq:dKdm}
\end{align}
Combining with \eqref{eq:dm-dtq} and \eqref{eq:dm-dk2},
\begin{equation}
  \frac{\partial \Lap K}{\partial \tq}
   = \frac{\partial \Lap K}{\partial m}\cdot
     \frac{s^{2}}{2m(\alpha+\kk s)} ,
  \qquad
  \frac{\partial \Lap K}{\partial \kk}
   = \frac{\partial \Lap K}{\partial m}\cdot
     \left(-\frac{ms}{2(\alpha+\kk s)}\right) .
  \label{eq:dK-dparams}
\end{equation}
These are inverted and convolved by exactly the same procedure as the forward
solution.
\end{derivationx}

\begin{checkbox}
Check 4c of the earlier forensic suite verifies \eqref{eq:dKdm} symbolically.
The assembled analytic Jacobian was compared against central differences of the
forward solver: maximum relative deviation $5.3\times10^{-7}$ \PREV. Because
$\alpha$ rescales the time axis and the pulse length rather than entering the
kernel as a parameter, its sensitivity is obtained by a single central
difference in $\alpha$; $\tq$ and $\kk$, whose identifiability is the subject
of this study, never use finite differences.
\end{checkbox}

\begin{derivationx}[Why $\Bnum=1$ makes $F$ singular, algebraically]
\label{der:why-singular}
This restates \cref{der:fisher} in the form used by the program. The two
columns of $\tilde J$ corresponding to $(\tq,\kk)$ are, on the ray,
$\tilde J_{:,\tq}=\tq J_{:,\tq}$ and $\tilde J_{:,\kk}=\kk J_{:,\kk}
=\alpha\tq J_{:,\kk}$. Using $J_{:,\tq}=-\alpha J_{:,\kk}$,
\begin{equation}
  \tilde J_{:,\tq} = -\alpha\tq\, J_{:,\kk} = -\,\tilde J_{:,\kk} .
  \label{eq:logJ-dependence}
\end{equation}
In log-parameters the two columns are exactly opposite, so
$\tilde J_{:,\tq}+\tilde J_{:,\kk}=\bm{0}$: the null direction in log
coordinates is $(1,1)$, i.e. the direction along which $\tq$ and $\kk$ are
scaled by the same factor, which is precisely motion along the ray.
\end{derivationx}

\begin{resultbox}
The algebraic proof of singularity precedes and is independent of any numerical
condition number. The condition numbers reported in \cref{sec:practical} are a
consequence, not the evidence.
\end{resultbox}

%=======================================================================
\section{Reparameterisation: what remains determined}
\label{sec:reparam}
%=======================================================================

\begin{derivationx}[Sensitivity in the $\bm{\phi}$ coordinates]
\label{der:reparam-sens}
Using $\bm{\phi}=(\alpha,\tq,\Bnum)$ and $\kk=\alpha\tq\Bnum$, the derivative
of the observable with respect to $\tq$ \emph{at fixed $\alpha$ and $\Bnum$} is
obtained by the chain rule, noting
$\partial\kk/\partial\tq|_{\alpha,\Bnum}=\alpha\Bnum$:
\begin{equation}
  \frac{\partial T}{\partial \tq}\bigg|_{\alpha,\Bnum}
   = \frac{\partial T}{\partial \tq}\bigg|_{\alpha,\kk}
   + \frac{\partial T}{\partial \kk}\cdot
     \frac{\partial \kk}{\partial \tq}\bigg|_{\alpha,\Bnum}
   = \frac{\partial T}{\partial \tq}
   + \alpha\Bnum\,\frac{\partial T}{\partial \kk} .
  \label{eq:reparam-chain}
\end{equation}
On the resonance ray $\Bnum=1$, and substituting \eqref{eq:antiparallel},
\begin{equation}
  \frac{\partial T}{\partial \tq}\bigg|_{\alpha,\Bnum=1}
   = -\alpha\frac{\partial T}{\partial \kk}
     + \alpha\cdot1\cdot\frac{\partial T}{\partial \kk}
   = 0 .
  \label{eq:reparam-degenerate}
\end{equation}
\end{derivationx}

\begin{checkbox}
Check 6b verifies $\partial_{\tq}\kk|_{\alpha,\Bnum}=\alpha\Bnum$ and that the
bracket in \eqref{eq:reparam-chain} vanishes on the ray. \NUM
\end{checkbox}

\begin{resultbox}[Localisation of the degeneracy]
In the coordinates $\bm{\phi}=(\alpha,\Bnum,\tq)$ the degeneracy is confined to
the single coordinate $\tq$: the derivative with respect to $\tq$ at fixed
$(\alpha,\Bnum)$ vanishes identically at $\Bnum=1$, while the derivatives with
respect to $\alpha$ and $\Bnum$ do not. The information carried by the data is
carried by $(\alpha,\Bnum)$; it is the resolution of that information into
$\tq$ and $\kk$ separately that fails.
\end{resultbox}

\subsection{Numerical evidence, kept separate from the theorem}
\label{sec:reparam-numeric}

The following is a numerical calculation and is \emph{not} part of the proof of
\cref{prop:main}. It quantifies how well $(\alpha,\Bnum)$ are determined for
one specific configuration.

\begin{calculationx}[Standard errors in both parameterisations]
\label{calc:reparam-num}
Configuration as in \cref{sec:practical}: $\tq=0.03$, $\tau_{\Delta}=0.04$,
$n=80$ points on $t\in[0.02,1.2]$, $\eta=2\%$, log-parameters, analytic
Jacobian. The standard error of $\Bnum$ is obtained from the delta method: with
$\ln\Bnum=\ln\kk-\ln\alpha-\ln\tq$ the gradient in log coordinates is
$g=(-1,-1,1)$ and
\begin{equation}
  \operatorname{s.e.}(\ln\Bnum)
   = \sqrt{g^{\!\top}\,\tilde F^{-1}\,g} .
  \label{eq:seB}
\end{equation}
\end{calculationx}

\begin{table}[H]
\centering
\caption{Standard errors in the two parameterisations. All values
         \NUM, script \texttt{validation/audit/reident.py} and
         \texttt{check\_numerical.py}. Configuration-specific.}
\label{tab:reparam}
\small
\begin{tabular}{@{}rrrrr@{}}
\toprule
$\Bnum$ & s.e.\ $\tq$ [\%] & s.e.\ $\kk$ [\%] &
s.e.\ $\alpha$ [\%] & s.e.\ $\Bnum$ [\%] \\
\midrule
0.50 & 12.37 & 21.56 & 0.94 & 9.56 \\
0.90 & 98.75 & 107.3 & 1.04 & 9.03 \\
0.98 & 526.6 & 535.0 & 1.06 & 8.90 \\
\textbf{1.00} & \textbf{divergent} & \textbf{divergent} &
  \textbf{0.94} & \textbf{4.00} \\
1.10 & 114.7 & 106.5 & 1.08 & 8.69 \\
1.28 & 45.64 & 37.80 & 1.11 & 8.39 \\
5.00 & 8.12  & 3.72  & 1.49 & 5.68 \\
\bottomrule
\end{tabular}
\end{table}

\begin{resultbox}
$\alpha$ is determined to about $1\%$ and $\Bnum$ to between $4\%$ and $10\%$
across the whole range \emph{including exactly at resonance}, while $\tq$ and
$\kk$ individually diverge. The data are not uninformative at $\Bnum=1$; the
split is what fails. This is a numerical result for the stated configuration
and does not extend the structural theorem.
\end{resultbox}

%=======================================================================
\section{Practical identifiability}
\label{sec:practical}
%=======================================================================

\begin{definitionx}[Practical identifiability]
\label{def:practical}
Recoverability of a parameter to a stated precision from finite, noisy data.
Unlike \cref{def:struct} it depends on the sampling design, the noise level and
the estimator, and it degrades continuously rather than switching on at a
point.
\end{definitionx}

\begin{definitionx}[Reference configuration]
\label{def:config}
Every number in this section and in \cref{sec:profile,sec:montecarlo,sec:bic}
refers to:
\begin{align}
  \text{true parameters:}&\quad \alpha=1,\ \ \tq=0.03,\ \
     \kk=\Bnum\,\alpha\,\tq ,
  \label{eq:cfg-1}\\
  \text{pulse:}&\quad \tau_{\Delta}=0.04 ,
  \label{eq:cfg-2}\\
  \text{sampling:}&\quad n=80 \ \text{points, uniform on}\
     t\in[0.02,\ 1.2] ,
  \label{eq:cfg-3}\\
  \text{noise:}&\quad \eta = 2\% \ \text{of the temperature range} ,
  \label{eq:cfg-4}\\
  \text{estimator:}&\quad \text{Levenberg--Marquardt on}\ \ln\bm{\theta}\
     \text{with analytic Jacobian} ,
  \label{eq:cfg-5}\\
  \text{working units:}&\quad L=1,\ t_{p}=0.04 .
  \label{eq:cfg-6}
\end{align}
\end{definitionx}

\begin{definitionx}[Objective function]
\label{def:objective}
\begin{equation}
  \mathrm{SSE}(\bm{\theta})
   = \sum_{i=1}^{n}
     \left[y_{i}^{\text{obs}} - T(1,t_{i};\bm{\theta})\right]^{2} ,
  \label{eq:sse}
\end{equation}
minimised over $\ln\bm{\theta}$. Under \eqref{eq:obsmodel} this is the maximum
likelihood estimator.
\end{definitionx}

\begin{calculationx}[Fisher standard errors versus $\Bnum$]
\label{calc:landscape}
For each $\Bnum$: build $\tilde J$ analytically, form
$\tilde F=\tilde J^{\!\top}\tilde J/\sigma^{2}$, invert, take
$\sqrt{\operatorname{diag}}$. Results in \cref{tab:landscape}.
\end{calculationx}

\begin{table}[H]
\centering
\caption{Practical identifiability versus $\Bnum$ for the configuration of
         \cref{def:config}. All values \NUM. Fisher from
         \texttt{reident.py}; Monte Carlo from \texttt{robust\_mc.py}
         (200 trials, spread reported as interquartile range divided by
         $1.349$, because the sampling distribution is heavy tailed near
         resonance).}
\label{tab:landscape}
\small
\begin{tabular}{@{}rrrrrrr@{}}
\toprule
$\Bnum$ & $\operatorname{cond}(\tilde F)$ &
$\operatorname{corr}(\tq,\kk)$ &
Fisher s.e.\ $\tq$ [\%] & MC s.e.\ $\tq$ [\%] &
MC $\tq$ 5--95\% & profile factor \\
\midrule
0.10 & $3.6\times10^{2}$ & $+0.767$ & 2.05  & ---   & ---          & --- \\
0.30 & $4.8\times10^{2}$ & $+0.926$ & 5.83  & ---   & ---          & --- \\
0.50 & $9.3\times10^{2}$ & $+0.974$ & 12.37 & 12.4  & 0.84--1.23   & 1.4 \\
0.70 & $2.7\times10^{3}$ & $+0.993$ & 27.06 & ---   & ---          & --- \\
0.90 & $2.5\times10^{4}$ & $+0.999$ & 98.75 & 103.9 & 0.09--5.88   & 35.5 \\
0.98 & $6.4\times10^{5}$ & $+1.000$ & 526.6 & ---   & ---          & --- \\
\textbf{1.00} & $\mathbf{1.4\times10^{16}}$ & \textbf{singular} &
  \textbf{divergent} & unbounded & unbounded & \textbf{unbounded} \\
1.02 & $6.5\times10^{5}$ & $+1.000$ & 542.5 & ---   & ---          & --- \\
1.10 & $2.7\times10^{4}$ & $+0.999$ & 114.7 & ---   & ---          & --- \\
1.28 & $3.5\times10^{3}$ & $+0.996$ & 45.64 & 48.9  & 0.35--1.85   & 3.5 \\
1.50 & $1.2\times10^{3}$ & $+0.987$ & 28.49 & ---   & ---          & --- \\
2.00 & $3.2\times10^{2}$ & $+0.957$ & 17.19 & ---   & ---          & --- \\
2.41 & $1.8\times10^{2}$ & $+0.927$ & 13.73 & 13.3  & 0.79--1.20   & --- \\
3.00 & $9.8\times10^{1}$ & $+0.882$ & 11.17 & ---   & ---          & --- \\
5.00 & $3.6\times10^{1}$ & $+0.757$ & 8.12  & 7.2   & 0.88--1.12   & 1.3 \\
\bottomrule
\end{tabular}
\end{table}

\begin{checkbox}
Checks 8-0.5, 8-1.28, 8-2.41, 8-5.0 reproduce the Fisher column to the quoted
precision; check 8e confirms $\operatorname{cond}(\tilde F)=1.4\times10^{16}$
at $\Bnum=1+10^{-7}$. \NUM
\end{checkbox}

\subsection{Derivation of the band boundaries}
\label{sec:band}

\begin{calculationx}[Threshold definition and root-finding]
\label{calc:band}
Define $\mathrm{se}_{\tq}(\Bnum)$ as the Fisher relative standard error of
$\tq$ from \cref{calc:landscape}. For a target $\Theta$ the band is the set
where the target is \emph{not} met:
\begin{equation}
  \mathcal{B}(\Theta)
   = \left\{\Bnum : \mathrm{se}_{\tq}(\Bnum) > \Theta\right\} .
  \label{eq:band-def}
\end{equation}
Since $\mathrm{se}_{\tq}$ increases monotonically towards $\Bnum=1$ from either
side, $\mathcal{B}(\Theta)$ is an interval $[\Bnum_{-},\Bnum_{+}]$ whose
endpoints solve $\mathrm{se}_{\tq}(\Bnum)=\Theta$. Each endpoint is located by
$50$ bisection steps, bracketing from $\Bnum=1$ outwards.
\end{calculationx}

\begin{resultbox}[Bands for the configuration of \cref{def:config}]
\begin{align}
  \mathrm{se}_{\tq} > 10\%  &:\quad
     \Bnum \in [0.441,\ 3.464] ,  &&\NUM
  \label{eq:band10}\\
  \mathrm{se}_{\tq} > 20\%  &:\quad
     \Bnum \in [0.628,\ 1.804] ,  &&\NUM
  \label{eq:band20}\\
  \mathrm{se}_{\tq} > 50\%  &:\quad
     \Bnum \in [0.817,\ 1.252] ,  &&\NUM
  \label{eq:band50}\\
  \mathrm{se}_{\tq} > 100\% &:\quad
     \Bnum \in [0.901,\ 1.116] .  &&\NUM
  \label{eq:band100}
\end{align}
Checks 8f and 8g reproduce the endpoints of \eqref{eq:band20} to
$\pm0.002$. The complementary design criterion, the region where the target
\emph{is} met, is $\Bnum\le\Bnum_{-}$ or $\Bnum\ge\Bnum_{+}$.
\end{resultbox}

\begin{checkbox}[Scope of these numbers --- mandatory qualification]
\eqref{eq:band10}--\eqref{eq:band100} are properties of the configuration of
\cref{def:config}. They are \emph{not} universal physical constants and would
move under a different sampling scheme, noise level or pulse length. They must
be recomputed for any materially different design; the program specification of
\cref{sec:program} therefore exposes $n$, $\eta$, $t_{\min}$, $t_{\max}$ and
$\tau_{\Delta}$ as inputs. Only the singularity at $\Bnum=1$ is
configuration-independent, because it is structural (\cref{prop:main}).
\end{checkbox}

%=======================================================================
\section{Profile likelihood}
\label{sec:profile}
%=======================================================================

\begin{definitionx}[Profile of the objective]
\label{def:profile}
Fix $\tq$ at a trial value $\tau$ and re-optimise the remaining parameters:
\begin{equation}
  \mathrm{SSE}_{\text{prof}}(\tau)
   = \min_{\alpha,\ \kk}\ \mathrm{SSE}(\alpha,\tau,\kk) .
  \label{eq:profile-def}
\end{equation}
The profile likelihood ratio statistic is
\begin{equation}
  \Delta(\tau)
   = \frac{\mathrm{SSE}_{\text{prof}}(\tau)
           - \min_{\tau'}\mathrm{SSE}_{\text{prof}}(\tau')}{\sigma^{2}} ,
  \label{eq:profile-stat}
\end{equation}
which under \eqref{eq:obsmodel} is asymptotically $\chi^{2}$ with one degree of
freedom. The $95\%$ confidence set is
\begin{equation}
  \left\{\tau : \Delta(\tau) \leq \chi^{2}_{1,0.95}\right\},
  \qquad \chi^{2}_{1,0.95} = 3.841 .
  \label{eq:profile-threshold}
\end{equation}
\end{definitionx}

\textbf{Why this is a stronger diagnostic than Fisher curvature.} The Fisher
standard error is the curvature of $\mathrm{SSE}$ at the optimum and therefore
assumes the objective is well approximated by a quadratic. Near a degeneracy
the objective develops a long flat valley and that approximation fails; the
Fisher error can then be finite while the true confidence set is unbounded.
The profile makes no quadratic assumption: it evaluates the objective along the
$\tq$ direction with all other parameters free, so it reports the actual extent
of the valley.

\begin{calculationx}[Procedure]
\label{calc:profile}
Grid: $53$ values of $\tau$ logarithmically spaced over
$\tq\times[10^{-1.3},10^{+1.3}]$, i.e. a factor of $\approx400$ either side.
At each node, \eqref{eq:profile-def} is solved by Levenberg--Marquardt on
$(\ln\alpha,\ln\kk)$ with two starting points to reduce the risk of a local
minimum. One noise realisation per $\Bnum$, seed fixed for reproducibility.
Script \texttt{validation/audit/run\_profile.py}.
\end{calculationx}

\begin{table}[H]
\centering
\caption{Profile-likelihood $95\%$ intervals for $\tq$. All values \NUM.
         Configuration of \cref{def:config}.}
\label{tab:profile}
\small
\begin{tabular}{@{}rlll@{}}
\toprule
$\Bnum$ & $95\%$ interval for $\tq$ &
As multiple of truth & Width factor \\
\midrule
0.50 & $[0.0267,\ 0.0378]$ & $[0.89,\ 1.26]$  & 1.4 \\
0.90 & $[0.0015,\ 0.0534]$ & $[0.05,\ 1.78]$  & 35.5 \\
\textbf{1.00} & $[0.0015,\ 0.5986]$ &
  $[0.05,\ 19.95]$ & \textbf{398, grid-limited} \\
1.28 & $[0.0238,\ 0.0846]$ & $[0.79,\ 2.82]$  & 3.5 \\
5.00 & $[0.0337,\ 0.0424]$ & $[1.12,\ 1.41]$  & 1.3 \\
\bottomrule
\end{tabular}
\end{table}

\begin{resultbox}
At $\Bnum=1$ the profile does not cross the threshold anywhere on a grid
spanning a factor of $400$: the interval is bounded only by the grid, so $\tq$
is \emph{unbounded} by the data, exactly as \cref{prop:main} requires. The
widths at the other $\Bnum$ values increase monotonically towards resonance and
are consistent with the Fisher column of \cref{tab:landscape} where the
quadratic approximation is valid.

The value at $\Bnum=1.28$, a factor of $3.5$, is the conditioning of a real
calibration (see \cref{sec:calibrations}) and brackets the factor $1.51$ to
$3.50$ by which two published evaluations of the same specimens disagree.
\end{resultbox}

%=======================================================================
\section{Monte Carlo parameter recovery}
\label{sec:montecarlo}
%=======================================================================

\begin{calculationx}[Procedure]
\label{calc:mc}
\begin{enumerate}[nosep,label=(\arabic*)]
  \item \textbf{Truth.} $\bm{\theta}^{\ast}=(1,\ 0.03,\ \Bnum\cdot0.03)$.
  \item \textbf{Noise generation.} Compute $y^{\ast}=T(1,t_{i};
        \bm{\theta}^{\ast})$; set $\sigma$ by \eqref{eq:sigma-def} with
        $\eta=0.02$; draw
        $y^{\text{obs}}=y^{\ast}+\mathcal{N}(0,\sigma^{2})$ independently per
        point. Generator: \texttt{numpy} PCG64, seeds $7000+i$ (fixed, so the
        study is exactly reproducible).
  \item \textbf{Trials.} $200$ per value of $\Bnum$.
  \item \textbf{Optimizer.} Levenberg--Marquardt (\texttt{scipy}
        \texttt{least\_squares}, \texttt{method='lm'}) on
        $\ln\bm{\theta}$, \texttt{xtol}=\texttt{ftol}$=10^{-13}$, with the
        analytic Jacobian of \cref{der:analytic-sens}.
  \item \textbf{Initial condition.} Started at the true parameters. This is the
        favourable case and is stated as such.
  \item \textbf{Bounds.} None; positivity is enforced by the log-parameterisation.
  \item \textbf{Summary statistic.} Because the sampling distribution is heavy
        tailed near resonance, the spread is reported as the interquartile
        range divided by $1.349$ (the Gaussian-consistent robust estimator),
        alongside the $5$--$95$ percentile range.
\end{enumerate}
Script \texttt{validation/audit/robust\_mc.py}.
\end{calculationx}

\begin{table}[H]
\centering
\caption{Monte Carlo recovery, $200$ trials. All values \NUM. Comparison
         column from \cref{tab:landscape}.}
\label{tab:mc}
\small
\begin{tabular}{@{}rllll@{}}
\toprule
$\Bnum$ & MC s.e.\ $\tq$ [\%] & Fisher s.e.\ $\tq$ [\%] &
$\alpha$ 5--95\% & $\Bnum$ 5--95\% \\
\midrule
0.50 & 12.4  & 12.37 & $[0.983,\ 1.017]$ & $[0.842,\ 1.149]$ \\
0.90 & 103.9 & 98.75 & $[0.982,\ 1.026]$ & $[0.064,\ 1.070]$ \\
1.00 & unbounded & divergent & unbounded & unbounded \\
1.28 & 48.9  & 45.64 & $[0.980,\ 1.019]$ & $[0.906,\ 1.410]$ \\
2.41 & 13.3  & 13.73 & $[0.978,\ 1.024]$ & $[0.905,\ 1.134]$ \\
5.00 & 7.2   & 8.12  & $[0.975,\ 1.030]$ & $[0.929,\ 1.108]$ \\
\bottomrule
\end{tabular}
\end{table}

\begin{resultbox}[Fisher and Monte Carlo agree]
Where the quadratic approximation is valid the two estimates agree closely:
$12.4$ versus $12.37$, $103.9$ versus $98.75$, $48.9$ versus $45.64$, $13.3$
versus $13.73$, $7.2$ versus $8.12$. They diverge only at $\Bnum=1$, where both
correctly report an unbounded result.
\end{resultbox}

\begin{checkbox}[A withdrawn claim, recorded so it is not reintroduced]
\label{check:withdrawn}
An earlier stage of this project reported a systematic factor-of-ten
discrepancy between Fisher and Monte Carlo uncertainties. That was traced to
finite-difference Jacobians contaminated by the solver noise floor and is
\textbf{withdrawn}. With the analytic Jacobians of \cref{der:analytic-sens} no
such discrepancy exists, as the table above shows. The withdrawn figure must
not appear as a scientific result anywhere in this project.

Similarly, an earlier optimizer-initialisation study reported a factor-of-23
spread in recovered $\tq$; it used the split-pulse inversion whose failure is
characterised in \cref{sec:numerics}, with the failure time inside the sampling
window. With the certified solver the effect is a factor of $4.7$ rising to
$9.3$ at a one-log-unit perturbation \PREV. The earlier figure is withdrawn.
\end{checkbox}

%=======================================================================
\section{Model selection by BIC}
\label{sec:bic}
%=======================================================================

\begin{definitionx}[BIC convention used]
\label{def:bic}
For $n$ observations, $k$ free parameters and Gaussian errors of unknown but
constant variance, the maximised log-likelihood satisfies
$-2\ln\mathcal{L}_{\max} = n\ln(\mathrm{SSE}/n) + \text{const}$, giving the
convention used throughout this work:
\begin{equation}
  \boxed{\ \mathrm{BIC} = n\ln\!\left(\frac{\mathrm{SSE}}{n}\right)
        + k\ln n\ }
  \label{eq:bic}
\end{equation}
Only differences $\Delta\mathrm{BIC}$ between models fitted to the \emph{same}
data are interpreted, so the omitted additive constant is irrelevant.
\end{definitionx}

\begin{checkbox}[Verification of the convention]
With $n=80$: $\ln n = \ln 80 = 4.382$. The penalty difference between GK
($k=3$) and Fourier ($k=1$) is $(3-1)\times4.382 = 8.764$. Hence at exact
resonance, where the two models fit identically well in expectation, BIC must
prefer Fourier by approximately $8.8$ units minus whatever the two extra
parameters recover by fitting noise. The observed value $6.64$ is consistent
with this and is \emph{not} an arbitrary number: it is the parsimony penalty
partially offset by noise-fitting. \DER
\end{checkbox}

\begin{calculationx}[Protocol]
\label{calc:bic}
Data generated by the full GK model at $\eta=2\%$, configuration of
\cref{def:config}. All four candidates of \cref{tab:models} fitted to the same
realisation with the same solver, same optimizer, same tolerances, and three
restarts to reduce local-minimum risk. Averaged over $15$ noise realisations.
$\Delta\mathrm{BIC}$ is measured from the best model in each row.
Script \texttt{validation/audit/run\_bic.py}.
\end{calculationx}

\begin{table}[H]
\centering
\caption{BIC model selection. All values \NUM. Lower is better; $0.00$ marks
         the selected model.}
\label{tab:bic}
\small
\begin{tabular}{@{}rrrrrl@{}}
\toprule
True $\Bnum$ & Fourier ($k{=}1$) & MCV ($k{=}2$) &
Nyíri ($k{=}2$) & GK ($k{=}3$) & Selected \\
\midrule
\textbf{1.00} & \textbf{0.00} & 236.58 & 4.03 & 6.64 &
  \textbf{Fourier} \\
1.28 & 23.14  & 253.21 & 1.38   & \textbf{0.00} & GK \\
2.41 & 166.73 & 302.12 & 44.21  & \textbf{0.00} & GK \\
5.00 & 267.00 & 349.25 & 105.29 & \textbf{0.00} & GK \\
\bottomrule
\end{tabular}
\end{table}

\begin{resultbox}
At $\Bnum=1$, BIC selects Fourier over GK from data generated by GK, with
$\Delta\mathrm{BIC}_{\text{GK}}=6.64$. This is the rank deficiency of
\cref{prop:main} recovered by a route that uses no derivative information and
no Fisher matrix. Away from resonance GK is selected decisively. MCV is
rejected everywhere by $237$ to $349$ units.
\end{resultbox}

\begin{checkbox}[Interpretation limit --- mandatory]
BIC measures statistical parsimony under the stated data-generating process,
noise model and observation design. It is used here as an identifiability
diagnostic. It is \emph{not} evidence that the material obeys Fourier's law:
at $\Bnum=1$ the two models are observationally identical, so the preference
reports what the data can distinguish, not the underlying physics.
\end{checkbox}

%=======================================================================
\section{Published calibration calculations}
\label{sec:calibrations}
%=======================================================================

Every value below is labelled. Values marked \LIT\ are transcribed from the
cited primary source; values marked \DER\ are computed here from those
literature values by \eqref{eq:Bdef}. No value is transferred between sources
without stating which source it came from.

\subsection{Worked arithmetic for each case}

\begin{calculationx}[$\Bnum$ from tabulated $(\alpha,\tq,\kk)$]
\label{calc:B-cases}
Wherever a source tabulates all three of $\alpha_{\mathrm{GK}}$, $\tq$ and
$\kk$, $\Bnum$ is obtained directly from \eqref{eq:Bdef}. The units cancel:
with $\alpha$ in $10^{-6}\,\si{\metre\squared\per\second}$, $\tq$ in
\si{\second} and $\kk$ in $10^{-6}\,\si{\metre\squared}$, the factors
$10^{-6}$ cancel between numerator and denominator, so the tabulated numbers
may be used directly. Explicitly, for the three Feh\'er--Kov\'acs basalt
specimens \cite{FeherKovacs2021}:
\begin{align}
  \Bnum(\SI{1.86}{\milli\metre})
   &= \frac{0.168}{0.61\times0.211}
    = \frac{0.168}{0.12871} = 1.3053 ,
  \label{eq:B-186}\\
  \Bnum(\SI{2.75}{\milli\metre})
   &= \frac{0.268}{0.61\times0.344}
    = \frac{0.268}{0.20984} = 1.2772 ,
  \label{eq:B-275}\\
  \Bnum(\SI{3.84}{\milli\metre})
   &= \frac{0.650}{0.68\times1.000}
    = \frac{0.650}{0.68} = 0.9559 .
  \label{eq:B-384}
\end{align}
The same source reports $1.295$, $1.272$ and $0.94$ respectively in its
Table~2 \LIT, so the largest disagreement is
$|0.9559-0.94| = 0.0159$.
\end{calculationx}

\begin{checkbox}
Check 11b: the maximum $|\Bnum_{\text{derived}}-\Bnum_{\text{reported}}|$ over
all cases where both exist is $0.0159$, occurring for the
\SI{3.84}{\milli\metre} specimen. Agreement to better than $0.02$ in every
case confirms that the present interpretation of the tabulated symbols matches
the authors' own. \NUM
\end{checkbox}

\begin{calculationx}[Both et al.\ capacitor, from the primary source]
\label{calc:both}
Both et al.\ \cite{Both2016} report, for the layered capacitor of thickness
\SI{3.9}{\milli\metre}, from a nonlinear regression on $2250$ data points:
\begin{align}
  \alpha &= (1.958 \pm 0.004)\times10^{-6}\ \si{\metre\squared\per\second}
    &&\LIT
  \label{eq:both-alpha}\\
  \tq &= (0.51 \pm 0.01)\ \si{\second} &&\LIT
  \label{eq:both-tau}\\
  \ell^{2} &= (1.53 \pm 0.03)\times10^{-6}\ \si{\metre\squared}
    &&\LIT
  \label{eq:both-l2}\\
  R^{2} &= 0.9996 &&\LIT
  \label{eq:both-R2}
\end{align}
With $\ell^{2}\equiv\kk$ (\cref{sec:notation-warning}),
\begin{equation}
  \Bnum = \frac{1.53\times10^{-6}}
               {1.958\times10^{-6}\times0.51}
        = \frac{1.53}{0.99858}
        = 1.5322 .
  \qquad\DER
  \label{eq:both-B}
\end{equation}
Propagating the three independent relative errors in quadrature,
\begin{align}
  \frac{\delta \Bnum}{\Bnum}
   &= \sqrt{\left(\frac{0.03}{1.53}\right)^{2}
          + \left(\frac{0.004}{1.958}\right)^{2}
          + \left(\frac{0.01}{0.51}\right)^{2}}
  \nonumber\\
   &= \sqrt{(0.01961)^{2}+(0.00204)^{2}+(0.01961)^{2}}
    = 0.02780 ,
  \label{eq:both-Berr}
\end{align}
so $\Bnum = 1.532 \pm 0.043$ \DER, a relative uncertainty of $2.78\%$.
\end{calculationx}

\begin{checkbox}
Checks 12c and 12d reproduce \eqref{eq:both-B} and \eqref{eq:both-Berr}. \NUM
\end{checkbox}

\begin{checkbox}[Which dataset is which --- mandatory distinction]
\label{check:both-vs-van}
V\'an et al.\ \cite{Van2017} also report a capacitor of thickness
\SI{3.9}{\milli\metre}, with $\alpha_{\mathrm{GK}}=2.13\times10^{-6}$,
$\tau=0.954\,\si{\second}$ and $b=2.23$ \LIT. These are \emph{different fitted
values} for the same specimen type, obtained by a different evaluation. They
are listed as a separate case in \cref{tab:calibrations} and are never merged
with \eqref{eq:both-alpha}--\eqref{eq:both-l2}. Only the Both et al.\ entry
carries reported uncertainties.
\end{checkbox}

\subsection{The V\'an et al.\ $\ell$/$b$ consistency question}
\label{sec:van-lb}

\begin{calculationx}[Reconciling the two tables of \cite{Van2017}]
\label{calc:van-lb}
That source defines \cite{Van2017}
\begin{equation}
  \hat\tau = \frac{\tau}{t_{p}},\quad
  \hat\alpha = \frac{\alpha t_{p}}{L^{2}},\quad
  \hat\ell = \frac{\ell}{L},\quad
  \hat b = \frac{\hat\ell^{2}}{\hat\tau\,\hat\alpha} = b ,
  \label{eq:van-defs}
\end{equation}
a different nondimensionalisation from \cref{def:scaling} (time scaled by
$t_{p}$, not $L^{2}/\alpha$). Substituting,
\begin{equation}
  b = \frac{(\ell/L)^{2}}{(\tau/t_{p})\,(\alpha t_{p}/L^{2})}
    = \frac{\ell^{2}}{L^{2}}\cdot\frac{t_{p}L^{2}}{\tau\,\alpha t_{p}}
    = \frac{\ell^{2}}{\tau\,\alpha} ,
  \label{eq:van-b}
\end{equation}
so $b$ is algebraically identical to $\Bnum$, with both $t_{p}$ and $L$
cancelling. Evaluating $\hat\ell^{2}/(\hat\tau\hat\alpha)$ from that source's
\emph{dimensionless} table alone gives $1.988$, $3.031$ and $1.782$ for
limestone, metal foam and leucocratic rock, against reported $2.17$, $3.04$ and
$1.77$ \DER: agreement to $0.01$ for the latter two.
\end{calculationx}

\begin{checkbox}[Root cause of the residual mismatch]
Combining that source's dimensionless table with its dimensional table does
\emph{not} reproduce $b$: for limestone the route
$\ell^{2}/(\tau\alpha_{\mathrm{GK}})$ gives $0.219$ against $2.17$. Testing the
internal consistency of the two tables through
$\hat\alpha \overset{?}{=} \alpha_{\mathrm{GK}}t_{p}/L^{2}$ with
$t_{p}=\SI{0.01}{\second}$:
\begin{center}
\begin{tabular}{@{}lrrr@{}}
\toprule
Sample & $\hat\alpha$ computed & $\hat\alpha$ reported & ratio \\
\midrule
Capacitor    & $1.400\times10^{-3}$  & $1.400\times10^{-3}$ & 1.000 \\
Limestone    & $1.021\times10^{-2}$  & $1.124\times10^{-3}$ & 9.08 \\
Metal foam   & $9.123\times10^{-4}$  & $9.120\times10^{-4}$ & 1.000 \\
Leucocratic  & $1.558\times10^{-2}$  & $1.560\times10^{-3}$ & 9.98 \\
\bottomrule
\end{tabular}
\end{center}
Two rows agree to four significant figures; two differ by factors of
$9.1$ and $10.0$ \DER. This is consistent with a units slip confined to two
rows of the published table.
\end{checkbox}

\begin{resultbox}
The authors' own reported $b$ values are used in \cref{tab:calibrations} for
cases 1--4, and the two affected rows are flagged. The discrepancy is
documented, not silently recomputed, because resolving it definitively would
require information not present in the published tables.
\end{resultbox}

\subsection{Master calibration table}

\begin{table}[H]
\centering
\caption{All twelve audited calibrations. R = reported by the source \LIT;
         C = calculated here \DER\ by \eqref{eq:Bdef}. Band is
         \eqref{eq:band20}, valid for the configuration of \cref{def:config}
         only. $\alpha$ in $10^{-6}\,\si{\metre\squared\per\second}$, $\tq$ in
         \si{\second}, $\kk$ in $10^{-6}\,\si{\metre\squared}$.}
\label{tab:calibrations}
\scriptsize
\begin{tabular}{@{}rlllllll@{}}
\toprule
\# & Source & Specimen & $\alpha$ & $\tq$ & $\kk$ & $\Bnum$ & In band? \\
\midrule
1 & V\'an+ \cite{Van2017} & capacitor \SI{3.9}{\milli\metre}
  & 2.13 R & 0.954 R & see \S\ref{sec:van-lb} & 2.23 R & no \\
2 & V\'an+ \cite{Van2017} & limestone \SI{1.7}{\milli\metre}
  & 2.950 R & 0.991 R & see \S\ref{sec:van-lb} & 2.17 R & no \\
3 & V\'an+ \cite{Van2017} & metal foam \SI{5.1}{\milli\metre}
  & 2.373 R & 0.402 R & see \S\ref{sec:van-lb} & 3.04 R & no \\
4 & V\'an+ \cite{Van2017} & leucocratic \SI{1.75}{\milli\metre}
  & 4.77 R & 1.32 R & see \S\ref{sec:van-lb} & 1.77 R & \textbf{yes} \\
\addlinespace
5 & earlier, in \cite{FeherKovacs2021} T1
  & basalt \SI{1.86}{\milli\metre}
  & 0.55 R & 0.738 R & 0.509 R & 1.243 R / 1.254 C & \textbf{yes} \\
6 & earlier, in \cite{FeherKovacs2021} T1
  & basalt \SI{2.75}{\milli\metre}
  & 0.604 R & 0.955 R & 0.670 R & 1.171 R / 1.162 C & \textbf{yes} \\
7 & earlier, in \cite{FeherKovacs2021} T1
  & basalt \SI{3.84}{\milli\metre}
  & 0.68 R & 0.664 R & 0.480 R & 1.060 R / 1.063 C & \textbf{yes} \\
\addlinespace
8 & Feh\'er--Kov. \cite{FeherKovacs2021} & basalt \SI{1.86}{\milli\metre}
  & 0.61 R & 0.211 R & 0.168 R & 1.295 R / 1.305 C & \textbf{yes} \\
9 & Feh\'er--Kov. \cite{FeherKovacs2021} & basalt \SI{2.75}{\milli\metre}
  & 0.61 R & 0.344 R & 0.268 R & 1.272 R / 1.277 C & \textbf{yes} \\
10 & Feh\'er--Kov. \cite{FeherKovacs2021} & basalt \SI{3.84}{\milli\metre}
  & 0.68 R & 1.000 R & 0.650 R & 0.940 R / 0.956 C & \textbf{yes} \\
11 & Feh\'er--Kov. \cite{FeherKovacs2021} & foam \SI{5.2}{\milli\metre}
  & 3.01 R & 0.304 R & 2.203 R & 2.395 R / 2.408 C & no \\
\addlinespace
12 & Both et al.\ \cite{Both2016} & capacitor \SI{3.9}{\milli\metre}
  & 1.958(4) R & 0.51(1) R & 1.53(3) R
  & 1.532(43) C & \textbf{yes} \\
\midrule
\multicolumn{8}{@{}p{0.97\textwidth}@{}}{\footnotesize
\textbf{Aggregate:} 8 of 12 inside the band \eqref{eq:band20}; 4 of 12 outside.
\textbf{Reported uncertainty:} case 12 only; the other eleven report none.
\textbf{Verification:} cases 5--7 verified against Table~1 of
\cite{FeherKovacs2021}, which tabulates the earlier evaluation alongside the
refined one, and Table~2 for their $\Bnum$; they were not traced further back
to the original measurement report. Cases 1--4 use the authors' reported $b$;
see \S\ref{sec:van-lb}.} \\
\bottomrule
\end{tabular}
\end{table}

\begin{checkbox}
Check 11a counts $8$ of $12$ inside $[0.628,1.804]$, matching the aggregate
above. The eight are cases 4--10 and 12. \NUM
\end{checkbox}

\subsection{Same-specimen discrepancy}
\label{sec:same-specimen}

\begin{table}[H]
\centering
\caption{Two independent evaluations of the same three basalt specimens, both
         tabulated in \cite{FeherKovacs2021}. All values \LIT.}
\label{tab:discrepancy}
\small
\begin{tabular}{@{}lrrrr@{}}
\toprule
Specimen & $\tq$ earlier [\si{\second}] & $\tq$ refined [\si{\second}] &
Ratio \DER & $R^{2}$ refined \\
\midrule
basalt \SI{1.86}{\milli\metre} & 0.738 & 0.211 & 3.50 & 0.9975 \\
basalt \SI{2.75}{\milli\metre} & 0.955 & 0.344 & 2.78 & 0.9964 \\
basalt \SI{3.84}{\milli\metre} & 0.664 & 1.000 & 1.51 & 0.9986 \\
\bottomrule
\end{tabular}
\end{table}

\begin{resultbox}
All six evaluations sit at $\Bnum=0.94$ to $1.31$. The profile-likelihood width
at $\Bnum=1.28$ is a factor of $3.5$ (\cref{tab:profile}), which brackets the
observed ratios $1.51$ to $3.50$. The spread is therefore consistent with and
quantitatively accounted for by the identifiability structure near
$\Bnum\approx1$.

\textbf{No causal claim is made.} The two evaluations differ in pulse
modelling, smoothing and fitting procedure as well as in conditioning, and the
published information does not permit these contributions to be separated.
\end{resultbox}

%=======================================================================
\section{The Feh\'er--Kov\'acs fitting domain}
\label{sec:domain}
%=======================================================================

\begin{resultbox}[Verbatim from the primary source]
Feh\'er and Kov\'acs state \cite{FeherKovacs2021}: \emph{``we still had to
restrict ourselves to a domain, which is $3 > \kappa^{2}/(\alpha\tau_{q})
\geq 1$. Its lower limit expresses the Fourier case.''} In the notation of this
document,
\begin{equation}
  1 \leq \Bnum < 3 .
  \label{eq:fitting-domain}
\end{equation}
\LIT
\end{resultbox}

\textbf{Significance of the lower boundary.} By \cref{prop:cancel}, $\Bnum=1$
is exactly the condition at which the GK response becomes the Fourier response,
and by \cref{prop:main} it is exactly where $(\tq,\kk)$ become structurally
non-identifiable. The accepted evaluation domain therefore has the degeneracy
as its lower boundary.

\begin{checkbox}[What is \emph{not} being asserted]
The restriction \eqref{eq:fitting-domain} is well motivated: room-temperature
measurements on these heterogeneous specimens indicate over-diffusive
behaviour, so the domain reflects the observed regime, and bounding a nonlinear
search improves its stability. \textbf{It is not claimed that the procedure of
\cite{FeherKovacs2021} is wrong.} Nor is it claimed that reported values
``happen to'' fall near $\Bnum=1$: the proximity is in part a consequence of
the admissible domain, and must be described as such. The same source
explicitly notes that GK theory itself is not restricted to this domain.
\end{checkbox}

%=======================================================================
\section{Calibration uncertainty consistency check}
\label{sec:uncertainty-check}
%=======================================================================

Case 12 is the only audited calibration reporting parameter uncertainties, so
it is the only opportunity to test the Fisher machinery of \cref{sec:fisher}
against a real published dataset.

\begin{calculationx}[Published relative uncertainty]
\label{calc:pub-unc}
From \eqref{eq:both-tau}:
\begin{equation}
  \frac{\delta\tq}{\tq}
   = \frac{0.01}{0.51}
   = 0.019608
   = 1.96\% .
  \qquad\text{PUBLISHED \LIT (arithmetic \DER)}
  \label{eq:pub-unc}
\end{equation}
\end{calculationx}

\begin{calculationx}[Predicted uncertainty under their design]
\label{calc:pred-unc}
\textbf{Assumptions, stated explicitly.} The Fisher calculation of
\cref{calc:landscape} is repeated with:
\begin{align}
  \Bnum &= 1.532 \quad\text{from \eqref{eq:both-B}} ,
  \label{eq:pred-1}\\
  n &= 2250 \quad\text{data points, as stated in \cite{Both2016}} \LIT ,
  \label{eq:pred-2}\\
  \eta &= 1\% \quad\text{assumed noise level (not stated in the source)} ,
  \label{eq:pred-3}\\
  &\text{all other settings as in \cref{def:config}} .
  \label{eq:pred-4}
\end{align}
\end{calculationx}

\begin{resultbox}
\begin{equation}
  \operatorname{s.e.}(\tq)\big|_{\text{their design}} = 2.55\% .
  \qquad\text{CALCULATED HERE \NUM}
  \label{eq:pred-unc}
\end{equation}
Check 12a, script \texttt{validation/calc/check\_numerical.py}. Compare
with the published $1.96\%$ of \eqref{eq:pub-unc}.
\end{resultbox}

\begin{checkbox}[Interpretation --- deliberately limited]
The two figures agree to within about $30\%$ of each other, from wholly
independent routes: one an asymptotic regression error reported by the original
authors, the other a Fisher prediction from the framework of this document.

\textbf{This is an external consistency check against a published real dataset,
not a validation.} The noise level \eqref{eq:pred-3} is assumed rather than
taken from the source, the sampling times are assumed uniform, and the reported
errors are themselves asymptotic and computed, as the authors state, assuming
exact $L$ and $t_{p}$. The check should be read as confirming that the
framework produces uncertainties of the right order for a real design, and as
demonstrating concretely why the band \eqref{eq:band20} is
configuration-dependent: the same $\Bnum$ with $2250$ low-noise points behaves
very differently from $80$ points at $2\%$ noise.
\end{checkbox}

%=======================================================================
\section{Validation anchor: reproduction of published results}
\label{sec:anchor}
%=======================================================================

\begin{definitionx}[Anchor]
\label{def:anchor}
R. Kov\'acs, \emph{Analytic solution of Guyer--Krumhansl equation for laser
flash experiments}, International Journal of Heat and Mass Transfer
\textbf{127}, Part A, 631--636 (2018), DOI
\texttt{10.1016/j.ijheatmasstransfer.2018.06.082} \cite{Kovacs2018}. The
version of record was confirmed on the publisher page; sections 4--5 are
paywalled and the equations used here were read from the open preprint
\texttt{arXiv:1804.05225}, whose equation numbering differs from the version of
record by one.
\end{definitionx}

\begin{calculationx}[Reproduction procedure]
\label{calc:repro}
\begin{enumerate}[nosep,label=(\arabic*)]
  \item Published figures extracted at native resolution
        ($4000\times1908$ px).
  \item Two independent digitisation passes: intensity-weighted centroid of the
        curve band, and envelope midpoint.
  \item Digitisation noise floor
        \begin{equation}
          \sigma_{d} = \sqrt{\overline{\left(\Delta_{AB}\right)^{2}}
                       + \operatorname{median}(\text{half-width})^{2}} ,
          \label{eq:sigmad}
        \end{equation}
        where $\Delta_{AB}$ is the pointwise difference between the passes.
  \item Published parameters used verbatim. \textbf{No fitting, no tuning, no
        adjustment of any kind.}
  \item Solver: the certified convolution solver of \cref{sec:numerics}.
  \item Metric: $\mathrm{NRMSE} = 100\times\mathrm{RMSE}/
        (\max T_{\text{pub}} - \min T_{\text{pub}})$.
\end{enumerate}
\end{calculationx}

\subsection{Truncated eigenfunction series required by Fig.~2}
\label{sec:series}

\textbf{Defect closed during Phase 3.} The convolution solver of
\cref{sec:numerics} returns the \emph{converged} solution. The anchor's Fig.~2
is a convergence study showing the solution \emph{truncated} at
$N=1,3,10,40$ eigenfunction modes, so reproducing it requires the truncated
series, which was not previously stated in this document. It is derived here.

\begin{derivationx}[Truncated series at the Fourier-resonance parameters]
\label{der:series}
Fig.~2 uses $\tq=\kk=0.02$, i.e. $\Bnum=1$, so by \eqref{eq:res-step5} the
operator is exactly the Fourier operator. Expand the dimensionless temperature
in the cosine basis $X_{n}(x)=\cos(n\pi x)$, which satisfies the insulated
rear-face condition and is the eigenbasis of $\partial_{xx}$ with eigenvalues
$-\beta_{n}$, $\beta_{n}=(n\pi)^{2}$:
\begin{equation}
  T(x,t)=a_{0}(t)+\sum_{n\geq1}a_{n}(t)\cos(n\pi x) .
  \label{eq:series-expansion}
\end{equation}
Projecting \eqref{eq:nd-energy} on mode $0$ gives
$\tau_{\Delta}\dot a_{0}=q_{0}(t)$, hence with $t_{c}=\min(t,\tau_{\Delta})$
\begin{equation}
  a_{0}(t)=\frac{1}{\tau_{\Delta}}\int_{0}^{t_{c}}q_{0}(u)\,\mathrm{d}u
          =\frac{1}{\tau_{\Delta}}
           \left[t_{c}-\frac{\sin(\omega t_{c})}{\omega}\right],
  \qquad \omega=\frac{2\pi}{\tau_{\Delta}} ,
  \label{eq:series-a0}
\end{equation}
which equals $1$ for $t\geq\tau_{\Delta}$, consistent with \eqref{eq:Tend}.
Mode $n\geq1$ obeys $\dot a_{n}+\beta_{n}a_{n}=A_{n}q_{0}(t)$ with
$A_{n}=2(-1)^{n}/\tau_{\Delta}$ at the rear face $x=1$, where
$\cos(n\pi)=(-1)^{n}$. Solving by the integrating factor and evaluating the
pulse integral only over its support gives
\begin{equation}
  a_{n}(t)=A_{n}\left[I_{1}(t)-I_{2}(t)\right],
  \label{eq:series-an}
\end{equation}
\begin{equation}
  I_{1}=\frac{\mathrm{e}^{-\beta_{n}(t-t_{c})}-\mathrm{e}^{-\beta_{n}t}}
             {\beta_{n}},
  \qquad
  I_{2}=\frac{\mathrm{e}^{-\beta_{n}(t-t_{c})}
              \left[\beta_{n}\cos(\omega t_{c})+\omega\sin(\omega t_{c})\right]
              -\beta_{n}\mathrm{e}^{-\beta_{n}t}}
             {\beta_{n}^{2}+\omega^{2}} .
  \label{eq:series-I}
\end{equation}
The rear-face solution truncated at $N$ modes is therefore
\begin{equation}
  \boxed{\ T_{N}(1,t)=a_{0}(t)+\sum_{n=1}^{N}a_{n}(t)\ }
  \label{eq:series-TN}
\end{equation}
\end{derivationx}

\begin{checkbox}
Script \texttt{validation/calc/derive\_series.py} verifies symbolically that
the closed form $I_{1}-I_{2}$ of \eqref{eq:series-I} equals the direct integral
$\mathrm{e}^{-\beta_{n}t}\int_{0}^{t_{c}}\mathrm{e}^{\beta_{n}u}q_{0}(u)
\mathrm{d}u$: residual exactly $0$. It also confirms
$a_{0}(\tau_{\Delta})=1$. \NUM
\end{checkbox}

\begin{checkbox}[Scope]
\eqref{eq:series-TN} is used \emph{only} to reproduce the anchor's truncation
study (test T4). Every other result in this document uses the converged
convolution solver of \cref{sec:numerics}. The two must agree as
$N\to\infty$, which is test T4b.
\end{checkbox}

\begin{table}[H]
\centering
\caption{Reproduction parameters, taken verbatim from \cite{Kovacs2018} \LIT.
         $\tau_{\Delta}=0.04$ in every case.}
\label{tab:anchor-params}
\small
\begin{tabular}{@{}llllr@{}}
\toprule
Anchor figure & Regime & $\tq$ & $\kk$ & Series terms / span \\
\midrule
Fig.~2 & convergence     & 0.02 & 0.02 & $N=1,3,10,40$; $t\in[0,1]$ \\
Fig.~3 & $\Bnum=1$ resonance & 0.02 & 0.02 & $N=40$; $t\in[0,1]$ \\
Fig.~5 & over-diffusive  & 0.02 & 0.20 & $N=10$; $t\in[0,2]$ \\
\bottomrule
\end{tabular}
\end{table}

\begin{resultbox}[Reproduction errors]
\begin{table}[H]
\centering
\small
\begin{tabular}{@{}lrrrrrr@{}}
\toprule
Figure & Points & RMSE & NRMSE & $\sigma_{d}$ &
$\mathrm{RMSE}/\sigma_{d}$ & within $\pm2\sigma_{d}$ \\
\midrule
Fig.~3 & 729 & 0.00384 & $\mathbf{0.384\%}$ & 0.00593 & 0.65 & 99.9\% \\
Fig.~5 & 743 & 0.00288 & $\mathbf{0.382\%}$ & 0.00603 & 0.48 & 99.5\% \\
Fig.~2 & 1999 & 0.00819 & $\mathbf{0.483\%}$ & --- & --- & --- \\
\bottomrule
\end{tabular}
\end{table}
All \NUM; checks 13-fig3 and 13-fig5 reproduce the first two rows exactly.
Fig.~2 is the pooled value over four truncation levels \PREV.
\end{resultbox}

\begin{checkbox}[Reading the ratio $\mathrm{RMSE}/\sigma_{d}<1$]
The reproduction error is \emph{below} the digitisation noise floor: the
residual is dominated by the precision with which the published curves can be
read, not by any discrepancy between the published solution and the present
implementation. A closer agreement could not be demonstrated from published
figures.

\textbf{Normalisation artefact, reported openly.} The $N=40$ curve of Fig.~2 is
nearly converged and spans a range of only $0.031$, so normalising its residual
by its own range gives $8.19\%$, while the same residual normalised by the full
range of the figure is $0.162\%$; the absolute error,
$2.7\times10^{-3}$, is comparable to the other curves \PREV.
\end{checkbox}

\begin{resultbox}[Scope of this validation --- mandatory]
This exercise validates the \emph{forward computational implementation} against
an independently authored published solution. It does \textbf{not} validate the
inverse problem, and no such claim is made. The inverse results of
\cref{sec:structural,sec:practical} rest on the analytic proof, the
arbitrary-precision test of \cref{sec:numdegen}, and the independent
statistical route of \cref{sec:bic}.

Distinguish carefully:
\begin{itemize}[nosep]
  \item \textbf{External validation:} agreement with results published by
        another author (this section).
  \item \textbf{Numerical cross-checking:} agreement between independent
        implementations of our own equations (\cref{sec:numerics}). This is
        not external validation.
\end{itemize}
\end{resultbox}

%=======================================================================
\section{Numerical certification}
\label{sec:numerics}
%=======================================================================

\subsection{Factorisation and why the pulse is not split}

\begin{derivationx}[Kernel/pulse factorisation]
\label{der:factorisation}
The rear-face transform \eqref{eq:That-rear} factorises as
\begin{equation}
  \Lap T(1,s) = \Lap K(s)\,\Lap Q_{0}(s) ,
  \qquad
  \Lap K(s) = \frac{m(s)}{\tau_{\Delta}\,s\,\sinh m(s)} ,
  \label{eq:factorisation}
\end{equation}
a material kernel times the pulse transform. $\Lap K$ has no poles on the
imaginary axis; $\Lap Q_{0}$ has the removable ones of
\cref{check:removable}. Therefore the kernel alone is inverted numerically and
the temperature is recovered by exact convolution:
\begin{equation}
  T(1,t) = \int_{0}^{\min(t,\tau_{\Delta})}
           k(t-u)\;q_{0}(u)\,\mathrm{d}u ,
  \qquad
  q_{0}(u) = 1-\cos\!\left(\frac{2\pi u}{\tau_{\Delta}}\right) ,
  \label{eq:convolution}
\end{equation}
$k$ being the inverse transform of $\Lap K$.
\end{derivationx}

\begin{derivationx}[The failure mode of the split formulation]
\label{der:pole-collision}
Suppose instead the pulse factor $1-\mathrm{e}^{-s\tau_{\Delta}}$ is handled by
the second-shift theorem and $\omega^{2}/[s(s^{2}+\omega^{2})]$ is inverted
alone. That factor retains genuine poles at $s=\pm\mathrm{i}\omega$. The fixed
Talbot rule evaluates the transform at
\begin{equation}
  s_{j} = r\,\theta_{j}\left(\cot\theta_{j} + \mathrm{i}\right),
  \qquad \theta_{j} = \frac{j\pi}{M},
  \qquad r = \frac{2M}{5t} .
  \label{eq:talbot-nodes}
\end{equation}
$\operatorname{Re}s_{j}=r\theta_{j}\cot\theta_{j}=0$ requires
$\cot\theta_{j}=0$, i.e. $\theta_{j}=\pi/2$, i.e. $j=M/2$, which is an integer
\emph{only when $M$ is even}. At that node
$s=\mathrm{i}\,r\pi/2$. This coincides with the pole $s=\mathrm{i}\omega$ when
\begin{equation}
  \frac{r\pi}{2} = \omega = \frac{2\pi}{\tau_{\Delta}}
  \quad\Longrightarrow\quad
  \frac{2M}{5t}\cdot\frac{\pi}{2} = \frac{2\pi}{\tau_{\Delta}}
  \quad\Longrightarrow\quad
  \boxed{\ t^{\ast} = \frac{M\,\tau_{\Delta}}{10}\ }
  \label{eq:tstar}
\end{equation}
\end{derivationx}

\begin{checkbox}[Prediction tested directly]
Check 7c verifies \eqref{eq:tstar} symbolically. Numerically, with
$\tau_{\Delta}=0.04$ \PREV:
\begin{center}
\begin{tabular}{@{}rrlll@{}}
\toprule
$M$ & $t^{\ast}$ & error at $0.98t^{\ast}$ & error at $t^{\ast}$ &
error at $1.02t^{\ast}$ \\
\midrule
20 (even) & 0.0800 & $1.9\times10^{-4}$ & $1.9\times10^{+11}$ &
  $1.0\times10^{-3}$ \\
40 (even) & 0.1600 & $3.9\times10^{-5}$ & $9.7\times10^{+10}$ &
  $7.2\times10^{-4}$ \\
41 (odd)  & 0.1640 & $1.1\times10^{-4}$ & $2.6\times10^{-4}$ &
  $2.1\times10^{-4}$ \\
60 (even) & 0.2400 & $6.8\times10^{-5}$ & $5.3\times10^{+9}$ &
  $5.7\times10^{-4}$ \\
64 (even) & 0.2560 & $9.2\times10^{-5}$ & $8.4\times10^{+10}$ &
  $5.8\times10^{-4}$ \\
\bottomrule
\end{tabular}
\end{center}
The error rises by ten to eleven orders of magnitude at exactly the predicted
time for every even $M$, and not at all for odd $M$.
\end{checkbox}

\begin{resultbox}[Corrected certification --- supersedes earlier statements]
\label{res:talbot}
\begin{enumerate}[nosep,label=(\arabic*)]
  \item The failure is a property of the numerical inversion, \emph{not} of the
        physics: the exact solution is regular at $t^{\ast}$.
  \item The remedy adopted is the convolution formulation
        \eqref{eq:convolution}, which inverts only the pole-free kernel, so the
        removable singularities never appear.
  \item Additionally $M$ is chosen \textbf{odd}, which removes the quadrature
        node from the imaginary axis entirely. Production value $M=41$.
  \item The naive remedy of increasing $M$ is actively harmful: by
        \eqref{eq:tstar} the failure time scales with $M$ and simply moves,
        while round-off grows.
  \item \textbf{An earlier certification described a ``safe window''
        $M\in[32,48]$ with instability at $M\approx56$--$68$. That description
        was obtained on a six-point probe set and does not survive testing on a
        dense grid; it is superseded by items (1)--(4) and must not be
        restored.}
  \item With the convolution formulation the accuracy plateau extends across
        $M\in[25,81]$, degrading beyond $M\approx91$ through round-off \PREV.
\end{enumerate}
\end{resultbox}

\subsection{Convergence, precision and cross-solver checks}

\begin{table}[H]
\centering
\caption{Numerical certification and error budget. Cross-solver agreement is
         numerical cross-checking, not external validation.}
\label{tab:certification}
\small
\begin{tabularx}{\textwidth}{@{}p{4.4cm}p{3.4cm}X@{}}
\toprule
Item & Result & Scope / label \\
\midrule
Convolution quadrature
  & $\leq10^{-13}$
  & 32 versus 64 Gauss--Legendre nodes \PREV \\
\addlinespace
Solver A vs.\ 50-digit reference
  & $10^{-11}$ to $10^{-13}$
  & probe times away from the steep rise \PREV \\
\addlinespace
Solver A vs.\ independent FD
  & $\approx10^{-7}$
  & full record; limited by the FD discretisation \PREV \\
\addlinespace
FD spatial convergence
  & observed order $p=2.00$;
    $2.6\times10^{-6}$ at $N_{x}=400$
  & Richardson reference \PREV \\
\addlinespace
FD temporal
  & $\approx10^{-9}$
  & relative tolerance $10^{-10}$ \PREV \\
\addlinespace
Arithmetic precision
  & $\leq8.9\times10^{-16}$
  & 15, 30 and 50 digits compared \PREV \\
\addlinespace
Fourier-resonance limit
  & $2.6\times10^{-6}$
  & $\kk=\alpha\tq$ against the analytical Fourier solution \PREV \\
\addlinespace
Fourier limit $\tq,\kk\to0$
  & $2.6\times10^{-6}$
  & against the analytical Fourier solution \PREV \\
\addlinespace
Non-negativity
  & $T\geq0$ in all cases tested
  & contrasts with negative-temperature artefacts reported for dual-phase-lag
    models \cite{Rukolaine2014} \PREV \\
\addlinespace
Steady state
  & $T(1,t\to\infty)=1.00000000$
  & confirms the normalisation \eqref{eq:Tend} \PREV \\
\addlinespace
\textbf{Total numerical error}
  & $\approx10^{-6}$
  & versus digitisation floor $6\times10^{-3}$ \\
\bottomrule
\end{tabularx}
\end{table}

\begin{resultbox}
The total numerical error of the framework is about $10^{-6}$, three orders of
magnitude below the digitisation noise floor of \cref{sec:anchor}. The
reproduction reported there is therefore digitisation-limited, not
solver-limited.
\end{resultbox}

%=======================================================================
\section{Master traceability table}
\label{sec:traceability}
%=======================================================================

\begin{longtable}{@{}p{3.5cm}p{1.6cm}p{2.6cm}p{1.5cm}p{1.5cm}p{1.6cm}@{}}
\caption{Traceability of every principal result. ``Program check'' marks the
         items the Phase-3 program must reproduce and the Phase-5 comparison
         must confirm.}
\label{tab:traceability}\\
\toprule
\textbf{Item} & \textbf{Eq./Tab.} & \textbf{Origin} &
\textbf{Derived here?} & \textbf{Num.\ checked?} & \textbf{Program check} \\
\midrule
\endfirsthead
\toprule
\textbf{Item} & \textbf{Eq./Tab.} & \textbf{Origin} &
\textbf{Derived here?} & \textbf{Num.\ checked?} & \textbf{Program check} \\
\midrule
\endhead
Energy balance & \eqref{eq:energy-dim} & literature
  \cite{Kovacs2018} & no & n/a & no \\
GK constitutive law & \eqref{eq:gk-dim} & literature
  \cite{GuyerKrumhansl1966a,GuyerKrumhansl1966b,Van2001} & no & n/a & no \\
Flux equation & \eqref{eq:gk-flux-dim} & derived & yes & n/a & no \\
Pulse BC & \eqref{eq:pulse-dim} & anchor \cite{Kovacs2018} & no & n/a & yes \\
Derived IC $\partial_t q(x,0)=0$ & \eqref{eq:ic-qdot} & derived & yes &
  n/a & yes \\
Scaling & \eqref{eq:scaling} & anchor \cite{Kovacs2018} & reproduced &
  n/a & yes \\
$\bar q_{0}=q_{\max}$ & \eqref{eq:qbar} & derived & yes & check 7d & yes \\
$T_{\mathrm{end}}$ & \eqref{eq:Tend} & anchor \cite{Kovacs2018} &
  reproduced & n/a & yes \\
Dimensionless system & \eqref{eq:nd-energy},\eqref{eq:nd-gk} & derived &
  yes & n/a & yes \\
$\hat\alpha\equiv1$ & \cref{der:nd-flux} & derived & yes & n/a & yes \\
Laplace elimination & \eqref{eq:lap-sub}--\eqref{eq:helmholtz-nd} & derived &
  yes & check 1a & no \\
\textbf{Propagation factor} & \eqref{eq:m2} & derived & yes &
  checks 1a,1b,2a--2c & yes \\
Spatial solution & \eqref{eq:qhat} & derived & yes & checks 5c,5d & yes \\
Rear-face transform & \eqref{eq:That-rear} & derived & yes & check 5d & yes \\
Pulse transform & \eqref{eq:Q0} & derived & yes & check 7a & yes \\
Removable poles & \eqref{eq:removable} & derived & yes & check 7b & yes \\
$\Bnum$ definition & \eqref{eq:Bdef} & adopted; ratio used by
  \cite{FeherKovacs2021,Van2017} & no & check 4a & yes \\
$\Bnum$ scale invariance & \eqref{eq:Binv} & derived & yes & check 4d & yes \\
$\kk=\alpha\tq\Bnum$ & \eqref{eq:kappa-from-B} & derived & yes &
  check 4b & yes \\
Param.\ Jacobian $=\alpha\tq$ & \eqref{eq:param-det} & derived & yes &
  check 4c & no \\
Fourier limit & \eqref{eq:lim-fourier} & derived & yes & check 3a & yes \\
MCV limit & \eqref{eq:lim-mcv} & derived & yes & check 3b & yes \\
MCV wave speed & \eqref{eq:mcv-speed} & derived & yes & check 3d & no \\
Nyíri limit & \eqref{eq:lim-nyiri} & derived; \cite{Nyiri1991} & yes &
  check 3c & yes \\
\textbf{Resonance cancellation} & \eqref{eq:res-step5} & derived; condition is
  prior art \cite{FulopKovacsVan2015,Kovacs2018,Fulop2018entropy} & yes &
  checks 5a,5b & yes \\
$\tq$-independence on ray & \eqref{eq:dTdtq-ray} & derived & yes &
  checks 5c,5d & yes \\
$\partial m/\partial\tq$ & \eqref{eq:dm-dtq} & derived & yes & check 5e & yes \\
$\partial m/\partial\kk$ & \eqref{eq:dm-dk2} & derived & yes & check 5f & yes \\
General ratio $-s/m^{2}$ & \eqref{eq:ratio-general} & derived & yes &
  check 5g & no \\
\textbf{Ratio $=-\alpha$} & \eqref{eq:ratio-resonance} & derived & yes &
  checks 5h,6a & yes \\
Anti-parallel sensitivities & \eqref{eq:antiparallel} & derived & yes &
  check 6a & yes \\
Rank one, null $(1,\alpha)$ & \eqref{eq:rank1},\eqref{eq:nullvec} & derived &
  yes & check 6b & yes \\
$\det F=0$ & \eqref{eq:detF} & derived & yes & check 6c & yes \\
$\operatorname{corr}=-1$ & \eqref{eq:corr} & derived & yes & check 6d & yes \\
$\partial_{m}\Lap K$ & \eqref{eq:dKdm} & derived & yes & verified & yes \\
Reparam.\ degeneracy & \eqref{eq:reparam-degenerate} & derived & yes &
  check 6b & yes \\
Numerical degeneracy & \eqref{eq:S-ray}--\eqref{eq:S-ratio} & computed &
  n/a & checks 7a--7c & yes \\
50-digit degeneracy & \eqref{eq:S-mp} & computed & n/a & check 7d & yes \\
Fisher landscape & \cref{tab:landscape} & computed & n/a & checks 8-* & yes \\
\textbf{Band $[0.628,1.804]$} & \eqref{eq:band20} & computed & n/a &
  checks 8f,8g & yes \\
Other bands & \eqref{eq:band10},\eqref{eq:band50},\eqref{eq:band100} &
  computed & n/a & same method & yes \\
Profile likelihood & \cref{tab:profile} & computed & n/a & rerun & yes \\
Monte Carlo & \cref{tab:mc} & computed & n/a & rerun & yes \\
BIC convention & \eqref{eq:bic} & standard & no & penalty check & yes \\
BIC results & \cref{tab:bic} & computed & n/a & rerun & yes \\
$\Bnum$ per calibration & \cref{tab:calibrations} & derived from
  \LIT & yes & check 11a,11b & no \\
Both et al.\ $\Bnum$ & \eqref{eq:both-B} & derived from \LIT & yes &
  check 12c & no \\
Both et al.\ $\delta\Bnum$ & \eqref{eq:both-Berr} & derived & yes &
  check 12d & no \\
Uncertainty cross-check & \eqref{eq:pred-unc} & computed & n/a &
  check 12a & yes \\
Fitting domain $1\le\Bnum<3$ & \eqref{eq:fitting-domain} & literature
  \cite{FeherKovacs2021} & no & n/a & no \\
Anchor NRMSE & \cref{sec:anchor} & computed & n/a & checks 13-* & yes \\
$t^{\ast}=M\tau_{\Delta}/10$ & \eqref{eq:tstar} & derived & yes &
  check 7c & yes \\
Error budget & \cref{tab:certification} & computed & n/a & \PREV & yes \\
\bottomrule
\end{longtable}

%=======================================================================
\section{Program implementation specification}
\label{sec:program}
%=======================================================================

This section specifies the Phase-3 program. \textbf{The program must implement
exactly the equations of this document and must contain no physics or equation
that does not appear above.}

\subsection{Module 1: forward solver}

\begin{description}[nosep,style=nextline]
\item[Inputs] $\alpha$ $[\si{\metre\squared\per\second}]$,
  $\tq$ $[\si{\second}]$, $\kk$ $[\si{\metre\squared}]$,
  $L$ $[\si{\metre}]$, $t_{p}$ $[\si{\second}]$,
  observation times $t_{i}$ $[\si{\second}]$.
\item[Internal scaling] $\hat t_{i}=\alpha t_{i}/L^{2}$,
  $\hat\tau_{q}=\alpha\tq/L^{2}$, $\hat\kappa^{2}=\kk/L^{2}$,
  $\hat\tau_{\Delta}=\alpha t_{p}/L^{2}$, and $\hat\alpha\equiv1$
  (\cref{der:nd-flux}). \textbf{The program must never mix
  $\hat\alpha\equiv1$ with the physical $\alpha$ in one expression};
  this specific error produced a $10^{6}$ mis-scaling in earlier work.
\item[Kernel] $\Lap K(s)=m/(\hat\tau_{\Delta}s\sinh m)$ with
  $m=\sqrt{s(1+\hat\tau_{q}s)/(1+\hat\kappa^{2}s)}$, principal branch
  (\cref{def:branch}).
\item[Inversion] fixed Talbot, nodes \eqref{eq:talbot-nodes},
  \textbf{$M$ odd, default $M=41$}. The program must reject even $M$ or emit a
  warning citing \eqref{eq:tstar}.
\item[Pulse handling] convolution \eqref{eq:convolution} with Gauss--Legendre,
  default $32$ nodes. \textbf{The second-shift/split formulation is
  prohibited} (\cref{res:talbot}).
\item[Overflow guards] evaluate $1/\sinh m$ as
  $2\mathrm{e}^{-m}/(1-\mathrm{e}^{-2m})$ and $\coth m$ as
  $(1+\mathrm{e}^{-2m})/(1-\mathrm{e}^{-2m})$; guard the removable $0/0$ when
  $\mathrm{e}^{-2m}$ underflows.
\item[Output] $T(1,t_{i})$, dimensionless, normalised so $T\to1$.
\end{description}

\subsection{Module 2: analytic Jacobian}

\begin{description}[nosep,style=nextline]
\item[Outputs] $\partial T/\partial\alpha$, $\partial T/\partial\tq$,
  $\partial T/\partial\kk$ at each $t_{i}$.
\item[Method] $\tq$ and $\kk$ analytically via \eqref{eq:dKdm} and
  \eqref{eq:dK-dparams}, inverted and convolved by the same route as the
  forward solution. $\alpha$ by a single central difference, since it rescales
  the time axis and pulse length rather than entering the kernel.
\item[Prohibited] finite differencing of the time-domain solver for $\tq$ or
  $\kk$ (\cref{check:withdrawn}).
\item[Acceptance] maximum relative deviation from Richardson-extrapolated
  central differences $<10^{-6}$.
\end{description}

\subsection{Module 3: independent cross-check solver}

Staggered-grid method of lines: $T$ at $N_{x}$ cell centres, $q$ at $N_{x}+1$
faces, so \eqref{eq:nd-bc} sits exactly on grid points; BDF integrator,
$N_{x}=400$, rtol $10^{-10}$, atol $10^{-12}$. Must share no code with
Module 1. Its agreement with Module 1 is numerical cross-checking, not
validation.

\subsection{Module 4: estimation and statistics}

\begin{description}[nosep,style=nextline]
\item[Estimator] Levenberg--Marquardt on $\ln\bm{\theta}$ with the Module-2
  Jacobian; xtol = ftol = $10^{-13}$.
\item[Fisher] $\tilde F=\tilde J^{\!\top}\tilde J/\sigma^{2}$,
  \eqref{eq:logJ},\eqref{eq:fisher-def}; report
  $\sqrt{\operatorname{diag}\tilde F^{-1}}$, $\operatorname{cond}\tilde F$,
  correlation, and s.e.$(\Bnum)$ by the delta method \eqref{eq:seB}.
\item[Profile likelihood] \eqref{eq:profile-def}--\eqref{eq:profile-threshold};
  $\geq53$ log-spaced nodes over at least $[10^{-1.3},10^{1.3}]\times\tq$,
  $\geq2$ restarts per node.
\item[Monte Carlo] $\geq200$ trials, fixed seeds, robust spread
  (IQR/1.349) reported alongside percentiles.
\item[BIC] \eqref{eq:bic} with $k$ from \cref{tab:models}; identical protocol
  across all four models.
\end{description}

\subsection{Required validation tests and expected values}

\begin{table}[H]
\centering
\caption{Acceptance tests the Phase-3 program must pass. Phase-5 compares
         program output against this column.}
\label{tab:acceptance}
\small
\begin{tabularx}{\textwidth}{@{}rXll@{}}
\toprule
\# & Test & Expected & Tolerance \\
\midrule
T1 & $m^{2}$ at $\kk=\alpha\tq$ equals $s/\alpha$ & exact &
  $10^{-14}$ \\
T2 & Anchor Fig.~3 NRMSE & $0.384\%$ & $\pm0.005$ \\
T3 & Anchor Fig.~5 NRMSE & $0.382\%$ & $\pm0.005$ \\
T4 & Anchor Fig.~2 pooled NRMSE, series \eqref{eq:series-TN} & $0.483\%$ & $\pm0.01$ \\
T4b & Series \eqref{eq:series-TN} at large $N$ vs convolution solver & agree & $10^{-4}$ \\
T5 & Spread on ray, $\tq$ over $100\times$, double precision &
  $\leq10^{-9}$ & --- \\
T6 & Control spread at $\Bnum=1.05$ & $\approx2.6\times10^{-2}$ &
  $\pm20\%$ \\
T7 & Fisher s.e.$(\tq)$ at $\Bnum=0.5$ & $12.37\%$ & $\pm0.1$ \\
T8 & Fisher s.e.$(\tq)$ at $\Bnum=1.28$ & $45.64\%$ & $\pm0.5$ \\
T9 & Band edges at $20\%$ & $0.628$, $1.804$ & $\pm0.005$ \\
T10 & $\operatorname{cond}\tilde F$ at $\Bnum=1+10^{-7}$ &
  $>10^{12}$ & --- \\
T11 & BIC at $\Bnum=1$: Fourier $0.00$, GK & $6.64$ & $\pm1.5$ \\
T12 & BIC at $\Bnum=2.41$ selects GK & GK & --- \\
T13 & Profile factor at $\Bnum=1.28$ & $3.5$ & $\pm0.5$ \\
T14 & Profile at $\Bnum=1$ exceeds grid & unbounded & --- \\
T15 & Monte Carlo vs Fisher at $\Bnum=0.5$ & agree within $20\%$ & --- \\
T16 & Even-$M$ split solver fails at $t^{\ast}=M\tau_{\Delta}/10$ &
  error $>10^{6}$ & --- \\
T17 & Odd-$M$ convolution solver at same $t$ & error $<10^{-3}$ & --- \\
T18 & FD vs Laplace solver, full record & $\leq10^{-6}$ & --- \\
T19 & Fourier limit $\tq,\kk\to0$ & $2.6\times10^{-6}$ & order of
  magnitude \\
T20 & $T(1,t\to\infty)$ & $1.00000000$ & $10^{-7}$ \\
T21 & $T\geq0$ everywhere & true & --- \\
T22 & $\Bnum$ identical across $L=1.86,2.75,3.84$ mm at fixed material &
  identical & $10^{-12}$ \\
T23 & $\Bnum$ for Both et al.\ values & $1.532$ & $\pm0.001$ \\
T24 & s.e.$(\tq)$ at $\Bnum=1.532$, $n=2250$, $\eta=1\%$ & $2.55\%$ &
  $\pm0.2$ \\
\bottomrule
\end{tabularx}
\end{table}

\subsection{Required figures}

F1 degeneracy on the ray with the $\Bnum=1.05$ control;
F2 anchor reproduction overlays with residuals against $\pm2\sigma_{d}$,
labelled as reproduction;
F3 Talbot pole-collision demonstration, even versus odd $M$;
F4 Fisher standard errors and $\operatorname{cond}\tilde F$ versus $\Bnum$
with the band shaded;
F5 profile likelihood curves with the $\chi^{2}_{1,0.95}$ threshold;
F6 BIC bar chart over the four models and four $\Bnum$ values;
F7 the twelve calibrations in $\Bnum$-space with the band and the domain
\eqref{eq:fitting-domain} shaded.

\subsection{Output format}

Machine-readable JSON for every table in this document, keyed by the equation
or table label, plus PDF figures. Phase-5 compares the JSON against
\cref{tab:acceptance} automatically.

%=======================================================================
\section{Final calculation audit}
\label{sec:audit}
%=======================================================================

\begin{table}[H]
\centering
\caption{Audit register. Full records in
         \texttt{validation/calc/checks\_symbolic.json} and
         \texttt{checks\_numerical.json}.}
\label{tab:audit}
\scriptsize
\begin{tabularx}{\textwidth}{@{}rp{3.1cm}p{3.5cm}Xl@{}}
\toprule
\# & Check & Input & Expected / Actual & Status \\
\midrule
1 & Symbolic algebra, $m^{2}$ & field equations & $s(1+\tq s)/(1+\kk s)$ /
  identical & PASS \\
2 & Dimensional analysis & all terms of
  \eqref{eq:energy-dim},\eqref{eq:gk-dim},\eqref{eq:m2} & consistent /
  consistent & PASS \\
3 & Limiting cases & $\tq\to0$, $\kk\to0$ & Fourier, MCV, Nyíri / all three
  recovered & PASS \\
4 & Parameter transformation & $\bm\theta\leftrightarrow\bm\phi$ &
  $\det=\alpha\tq$ / $\alpha\tq$ & PASS \\
5 & Structural identifiability & $\kk=\alpha\tq$ &
  $\partial_{\tq}T=0$, ratio $=-\alpha$ / both exact & PASS \\
6 & Fisher rank & $J$ on the ray & rank 1, $\det F=0$, corr $-1$ / all three
  & PASS \\
7 & Numerical degeneracy & $\tq$ over $100\times$ &
  $\leq10^{-9}$ / $1.26\times10^{-11}$; 50-digit $7.7\times10^{-30}$ &
  PASS \\
8 & Profile likelihood & 53-node grid & monotone widening; unbounded at
  $\Bnum=1$ / as expected & PASS \\
9 & Monte Carlo & 200 trials & agrees with Fisher / $12.4$ vs $12.37$ etc.
  & PASS \\
10 & BIC & 4 models, 15 realisations & Fourier at $\Bnum=1$ /
  $0.00$ vs GK $6.64$ & PASS \\
11 & Calibration arithmetic & 12 cases & 8 inside; max dev $0.016$ /
  8; $0.0159$ & PASS \\
12 & Both et al.\ uncertainty & their design & $\approx2.5\%$ vs published
  $1.96\%$ / $2.55\%$ & PASS \\
13 & Anchor reproduction & published params & $0.384\%$, $0.382\%$ /
  identical & PASS \\
14 & Stale-number scan & superseded values & absent / absent & PASS \\
15 & Citation--source check & 14 entries & all resolve / all resolve &
  PASS \\
16 & Equation--source check & all numbered equations & origin assigned /
  assigned & PASS \\
\bottomrule
\end{tabularx}
\end{table}

\subsection{Discrepancies found and resolved during this audit}

\begin{enumerate}[label=(D\arabic*)]
\item \textbf{Sympy radical reduction.} Checks 6a, 6b and 6d initially failed
because the algebra system does not reduce
$\sqrt{s^{2}\tq^{2}+2s\tq+1}$ to $(1+s\tq)$. The polynomial factors exactly as
$(1+s\tq)^{2}$ and both factors are positive for $s,\tq>0$, so the reduction is
valid; numerical spot checks confirm the difference is zero to machine
precision. The check script now applies the reduction explicitly and all three
pass. \textbf{This was a limitation of the tool, not an error in the
mathematics.}

\item \textbf{Tuple comparison in the audit recorder.} The dimensional checks
initially raised a type error because dimension pairs were compared with
symbolic subtraction. Fixed in the recorder; no mathematical content affected.

\item \textbf{Superseded Talbot certification.} An earlier ``safe window
$M\in[32,48]$, unstable $M\approx56$--$68$'' was obtained on a six-point probe
set and fails on a dense grid. Replaced by the derived pole-collision mechanism
\eqref{eq:tstar} and the convolution remedy (\cref{res:talbot}).

\item \textbf{Withdrawn Fisher/Monte-Carlo discrepancy and optimizer spread.}
Both traced to numerical artefacts and withdrawn; recorded in
\cref{check:withdrawn} so they are not reintroduced.
\end{enumerate}

\subsection{Stale-number scan}

The following superseded values were searched for and are \textbf{absent} from
this document except where explicitly identified as withdrawn:
$0.661$, $1.723$ (old band); ``6 of 11'', ``7 of 11'' (old aggregates);
factor $23$ (old optimizer spread); tenfold Fisher/Monte-Carlo discrepancy;
$\Bnum\leq0.20$ (old design branch); $M\in[32,48]$ as a safe window.

\subsection{Remaining blockers}

\begin{enumerate}[label=(B\arabic*)]
\item \textbf{Authenticated citation-index search not performed.} Scopus and
Web of Science are authentication-walled in the working environment. This
affects only statements about the prior literature, and this document makes
none: it records calculations, not novelty claims. It remains a blocker for the
Phase-6 manuscript.
\item \textbf{Anchor sections 4--5 paywalled.} Equations were read from the
open preprint \texttt{arXiv:1804.05225}; the bibliographic record was confirmed
against the publisher page. Equation numbers of the anchor are therefore never
cited.
\item \textbf{Raw heat-pulse time series unavailable.} The calibration audit of
\cref{sec:calibrations} necessarily works from tabulated fitted values. A
direct re-inversion of the original records is not possible with published
material.
\item \textbf{Noise level in \cref{calc:pred-unc} is assumed}, not taken from
\cite{Both2016}, which does not state it. The check is labelled a consistency
check for this reason.
\end{enumerate}

\bibliographystyle{unsrt}
\bibliography{references}

\end{document}
